\documentclass[leqno,12pt]{amsart}
% \documentclass[leqno]{amsart} AMS packages:
\usepackage{amscd} \usepackage{graphicx}
\usepackage{tikz}
\usepackage{tikz-cd}
\usepackage[margin=25mm]{geometry}
\usepackage{booktabs}
% Author's extra math macros and changes to existing macros:
\usepackage{blaom} \usepackage{hyperref}
% to increase height of rows in tables:
\renewcommand{\arraystretch}{1.2}
% Uncomment the following two lines to display reference labels:
%\usepackage[left]{showlabels}
%\renewcommand{\showlabelfont}{\footnotesize\color{blue}}
% Macros just for this paper:
  \newcommand{\varcheck}[1]{{#1^\wedge}}
  \newcommand{\model}{\automorphism({TM};A)}
%  \newcommand{\ii}{\mathrm{I\!I}}
  % \renewcommand{\contentsname}{\vspace{-\baselineskip}}
  % Drafting options: \usepackage{my_page_layout}
  % \usepackage{double_spacing}
% \usepackage{pdfsync}
%
\begin{document}
\bibliographystyle{plain}
%
\title%
{Filtered Lie algebroids and parabolic geometry} \author{Anthony D.~Blaom}
%\mbox{}\\  \\ {\sffamily --- Notes on work in progress ---}}
\address{E-mail: {\tt anthony.blaom@gmail.com}}%
% \keywords{Lie algebroid, Cartan algebroid, equivalence, geometric
% structure, Cartan connection, deformation, pseudogroup,
% prolongation}%
% \subjclass{Primary
% 53C15, % General geometric structures on manifolds
% 58H15% Deformations of structures
% ; Secondary 53B15, % Other connections
% 53C07, % Special connections and metrics on vector bundles
% 53C05, % Connections, general theory
% 58H05, % Pseudogroups and differential groupoids
% 53A55, % Differential invariants (local theory), geometric objects
% 53A30, % Conformal differential geometry
% 58A15% Exterior differential systems (Cartan theory)
% }
\thispagestyle{empty}
\begin{abstract}
  The filtered Lie algebroids introduced here generalize the filtered manifolds in Tohru
  Morimoto's study of finite-type geometric structures. We describe a development of
  parabolic geometry using Lie groupoids whose associated Lie algebroids are
  filtered.
\end{abstract}
\maketitle

%\tableofcontents
%%%%%%%%%%%%%%%%%%%%%%%%%%%%%%%%% Body starts
%%%%%%%%%%%%%%%%%%%%%%%%%%%%%%%%% here %%%%%%%%%%%%%%%%%%%%%%%
% \noindent%
Work in progress. Not for general circulation.

\section{Filtered Lie algebroids}\label{introduction}

\subsection{Definition}\label{filtered}
In this article, a filtration
$$\cdots A^{-1}\supset A^0 \supset A^1 \supset A^2 \supset\cdots $$ of some object $A$
will always be understood to satisfy $\cap_j A^j =0$ and $\cup_j A^j = A$. All filtrations
are descending. If $A$ is a vector bundle, then we call the filtration a {\df Lie
  algebroid filtration} if $A$ is a Lie algebroid and:

\begin{conditions}
\item[\bf F1.] The induced filtration
  $\cdots \Gamma(A^{-1})\supset \Gamma(A^0) \supset \Gamma(A^1) \supset \Gamma(A^2)
  \supset\cdots $ of the section space $\Gamma(A)$ is a filtration of Lie algebras,
  meaning $[X, Y] \in \Gamma(A^{j+k})$ for $X \in \Gamma(A^j)$ and $Y \in \Gamma(A^k)$.
\item[\bf F2.] The kernel of the anchor contains $A^0$.\label{fa2}
\end{conditions}

A filtered Lie algebroid $A$ is {\df firmly filtered} if $A^0$ and the filtration of $A^0$
are invariant under the canonical representation of $A$ on its anchor kernel, i.e., if
$\Gamma(A^j) \subset \Gamma(A)$ is an ideal for all $j \ge 0$.

\begin{proposition} Let $A$ be a filtered Lie alebroid. Then:
  \begin{conditions}
  \item The associated graded vector
    bundle $\graded(A)$ is a graded Lie algebra bundle. \label{b1}
  \item If $A$ is firmly filtered, then $\graded(A)$ is the Lie algebra bundle direct sum
    of its negative and non-negative parts.\label{rt2}
    \end{conditions}
  \end{proposition}
\noindent Condition \eqref{rt2} means that $\oplus_{i<0}\graded(A)_i$ and
$\oplus_{i\ge 0}\graded(A)_i$ are fibre-wise ideals of $\graded(A)$.

\begin{proof}
  Let $X \in \Gamma(A^j) $ and $Y \in \Gamma(A^k)$ be given. To prove \eqref{b1}, it
  suffices to show that
  $$ [fX, Y] = f[X, Y] \mod \Gamma(A^{j+k+1})$$
  for any smooth function $f$ on the base manifold. If $\#$ denotes
  the anchor, then the necessary and sufficient condition is
  that
  $$df(\#Y)X \in \Gamma(A^{j + k + 1}).$$
  If $k \ge 0$ then this is holds by F2. If instead $k \le -1 $, then $j + k + 1 \le j$,
  implying $A^j \subset A^{j + k + 1}$ and the condition is still satisfied.

  Suppose $A$ is firmly filtered. Then it is clear that $\oplus_{i<0}\graded(A)_i$ and
  $\oplus_{i\ge 0}\graded(A)_i$ are subobjects in the category of Lie algebra bundles. To
  prove \eqref{rt2}, it suffices to show that sections of one subbundle commute with
  sections of the other. To this end, suppose $X \in \Gamma(A^k)$ and $Y \in \Gamma(A^j)$,
  with $k<0$ and $j\ge 0$. Then $j + k + 1 \le j$, which implies $A^j \subset
  A^{j+k+1}$. Since $A$ is firmly filtered, $\Gamma(A^j) \subset \Gamma(A)$ is an ideal,
  so that $[X, Y] \in \Gamma(A^j) \subset \Gamma(A^{j + k + 1})$. It follows that
  $[\graded_j(X), \graded_k(Y)]=0$.
\end{proof}

The remainder of the section is devoted to examples of filtered Lie algebroids. All but
the first class of examples are always firm. We include some observations foreshadowing
general featutes of parabolic geometry to be presented later.

\subsection{ Filtered Lie algebras and Lie algebra bundles} Evidently, every filtered
Lie algebra, and more generally every filtered Lie algebra bundle, is a filtered Lie
algebroid.

\subsection{ Filtered manifolds} Axiom F2 forces Lie algebroid filtrations on a tangent
bundle $TM$ to be concentrated in negative degrees ($T^jM =0$ for $j\ge 0$). These
filtrations are trivially firm. A smooth manifold $M$ whose tangent bundle is so
structured is called a {\df filtered manifold}, a notion introduced and applied to
geometric structures by Morimoto \cite{Morimoto_93}. The anchor
$\# \colon A \rightarrow M$ of any {\em transitive} filtered Lie algebroid $A$ pushes the
filtration down to a filtration of $TM$. That is, the base of any transitive filtered Lie
algebroid is a filtered manifold.\footnote{One can drop the requirement of transitivity,
  but only at the expense of dropping our requirement $\cap_jA_j=0$ in the definition of a
  filtration.} Conversely, every transitive Lie algebroid $A$ over a filtered manifold $M$
gets a filtration by pulling back the filtration of $TM$, although this filtration is
stable in positive degrees: $A^j = A^0$ for $j\ge0$.\label{yt}

\subsection{Trivial filtrations} Suppose $A$ is any Lie algebroid whose anchor kernel has
constant rank. Choosing $A^0$ to be this kernel, and defining $A^j=A$ for $j<0$ and
$A^j = 0$ for $j>0$, we get a firm Lie algebroid filtration that we call the {\df trivial
  filtration} of $A$.

\subsection{ Filtered Klein geometries}\label{eg2}
Let $G$ be a Lie group with Lie algebra
${\mathfrak g}$ and $P \subset G$ a subgroup. Furthermore, suppose that ${\mathfrak g} $
admits a $P$-invariant filtration such that ${\mathfrak g}^0$ is the Lie algebra of
$P$. In particular, as we recall later, a filtration of this kind exists whenever $G/P$ is
a generalized flag manifold.  Write $M = P/G$ and let $A={\mathfrak g} \times M$ be the
action algebroid. Then we obtain a filtration of $A$, as a vector bundle, by defining
$A^j(m) = \Adjoint_g({\mathfrak g}^j)\times \{m\}$, where $g \in G$ is any element
satisfying $m = g \cdot m_0$. Here $m_0$ is the left coset of $P $ containing the identity
of $G$. It is too not difficult to see that $A$ is a firmly filtered Lie algebroid
(Proposition \ref{homogeneous}).

Significantly, the corresponding graded Lie algebra bundle is {\em not}
modelled on $\graded({\mathfrak g})$ but modelled on the ordinary direct sum of the graded
Lie algebras $\oplus_{i<0}\graded({\mathfrak g})_i$, equipped with {\em negative} bracket,
and $\oplus_{i\ge 0}\graded({\mathfrak g})_i$. So $\graded(A)$ has {\em two} different,
although related, Lie algebra bundle structures --- the one inherited from the filtered Lie
algebroid, and the ``flat'' one defined by ${\mathfrak g} $. On $A^0$ they coincide.

Note that the filtration of $A$ is firm, irrespective of whether
${\mathfrak g}_j \subset {\mathfrak g} $ is an ideal for any $j$.

\subsection{ Projective structures.}\label{projective} Let $A \subset J^2(TM)$ be the abstract PDE you solve
to find the infinitesimal isometries of some projective structure on $M$. Then $A$ is a
subalgebroid mapped isomorphically onto $J^1(TM)$ by the projection
$J^2(TM) \rightarrow J^1(TM)$, whose kernel intersects $A$ in the image $A^1$ of the
natural injection $j \colon T^*\!M \rightarrow S^2(T^*\!M)\otimes TM$ defined by
$$j(\alpha)(V_1, V_2) = \alpha(V_1)V_2 + \alpha(V_2)V_2.$$ See \cite[\S5.9]{Blaom_12}. If we define
$A^0$ to be the anchor kernel of $A$, then $A^1\subset A^0 \subset A^{-1}=A$ is a firm Lie
algebroid filtration. The corresponding graded Lie algebra bundle is
$$\graded(A) = TM \oplus \gl(TM) \oplus T^*\!M,$$ where $\gl(TM) \oplus T^*\!M $ is the semidirect product, and $TM$ and $\gl(TM) \oplus T^*\!M $
are ideals.

While not immediately obvious, $\graded(A)$ has a second, natural, {\em semisimple,} Lie
algebra bundle structure, modeled on $\speciallinear({\mathbb R}, n + 1)$,
$n= \dimension(M)$, and compatible with the grading. The bracket for this structure, which
we will denote $T$, extends the existing bracket on $\graded(A^0)=\gl(TM) \oplus T^*\!M $
with the new relations
\begin{equation}
  T(\phi, V) = \phi V, \quad
    T(V, \alpha) = j(\alpha)V, \label{projective_bracket}
\end{equation}
for sections $\phi$, $V$ and $\alpha $ of $\gl(TM)$, $TM$ and $T^*\!M $ respectively.
Significantly, this bracket coincides with the flat one referenced in \ref{eg2}, when $M$
is the projective sphere. \label{zt}

\subsection{Conformal structures}\label{conformal}
Let $B \subset J^1(TM)$ be the abstract PDE you solve to find the infinitesimal isometries
of a conformal structure on $M$ \cite[Section 12]{Blaom_12}. Then $B$ is a transitive
subalgebroid whose anchor kernel $\co(TM)\subset \gl(TM)$ consists of the infinitesimal
conformal automorphisms of tangent spaces. Now $B$ acts on $TM$ via restricted adjoint
action, and so acts on $T^*\!M $. Regarding $T^*\!M $ as a commutative Lie algebra bundle,
we construct the semidirect product $A = B \oplus_{\adjoint} T^*\!M $ (what we shall come to call
the {\df algebraic prolongation} of $B$). Then, if we put $A^1 = T^*\!M $,
$A^0 = TM \oplus \co(TM)$ (the anchor kernel of $A$), and $A^{-1}=A$, then $A$ becomes a
firmly filtered Lie algebroid. The corresponding graded Lie algebra bundle is
$$\graded(A)=TM \oplus \co(TM) \oplus T^*\!M,$$ where $\co(TM) \oplus T^*\!M $ is the
semidirect product, and $TM$ and $\co(TM) \oplus T^*\!M $ are ideals.

Again, there is a semisimple Lie algebra bundle structure on
$\graded(A)$ compatible with the grading, in this case
modeled on $\so({\mathbb R}, n + 1, 1)$, and coinciding with the flat one of \ref{eg2}
when $M$ is the conformal sphere. After redefining the injection $j$ above,
\begin{equation*}
    j(V_1, V_2) = \alpha(V_1) V_2 + \alpha(V_2)V_2 - \sigma(V_1,V_2)\sigma^{-1}(\alpha),
\end{equation*}
now a map
$j \colon T^*\!M\rightarrow (S^2(T^*\!M)\otimes TM) \cap (T^*\!M \otimes \co(TM)) $, the
semisimple bracket $T$ extends the existing bracket on
$\graded(A^0)=\co(TM) \oplus T^*\!M$, with the same additional relations
\eqref{projective_bracket} above.  Here $\sigma $ is any metric in the conformal class.

\section{Semidirect product structure}\label{semidirect}

Nowhere in this section are Lie algebroids or Lie groupoids required to be transitive.

\subsection{The reduction of firmly filtered Lie algebroids}\label{reduction}
If a filtration of a Lie algebroid $A$ is firm, then the canonical action $\adjoint $ of
$A$ on its anchor kernel preserves $A^0$ and so descends to an action of $A$ on
$\graded(A^0)$, also denoted $\adjoint$, preserving the grading, and the bracket on
$\graded(A^0)$. Moreover, fixing some $n \ge 1$ (but typically $n=1$) it follows from the
filtration property that $\adjoint$ restricts to a trivial action of the subalgebroid
$A^n$ on $\graded(A^0)$, delivering an action of $A/A^n$ (a Lie algebroid, because
$\Gamma(A^n) \subset \Gamma(A)$ is an ideal) on $\graded(A^0)$ that preserves the grading
and bracket. We call $A/A^n$ the {\df $n$th-order reduction} of $A$, or the {\df Levi
  reduction} of $A$, in the special case $n=1$. If $A$ is transitive and $A^0$ is the
anchor kernel, then the reductions are also transitive.

Now $B=A/A^n$ is itself filtered, with $B^i=0$ for $i \ge n$. Its representation
$E=\graded(A^n) \subset \graded(A^0)$ is canonically filtered, with $E^i=E$ for $i \le n$.
It is not hard to see that the Lie algebroid semidirect product $B \oplus_{\adjoint} E$ is
accordingly filtered, with $(B \ltimes E)^i=B^i \oplus E^i$, $i \in {\mathbb Z} $.  The
present section answers the question, ``When can $A$ be recovered as the Lie algebroid
semidirect product $B \ltimes E$?'' and formulates an analagous global result for Lie
groupoids whose associated Lie algebroids are filtered.

\subsection{Splittings of a filtered Lie algebra}
\label{splittings}
Recall that an element $E$ of a graded Lie algebra
${\mathfrak g}=\oplus_j {\mathfrak g}_j $ is a {\df grading element} if $[E , X] = j X$,
whenever $X \in {\mathfrak g}_j$. Necessarily, $E \in {\mathfrak g}_0$. If ${\mathfrak g}$
is merely filtered, then the existence of a grading element $E \in \graded({\mathfrak g})$
has strong consequences for the Lie algebra. For, choosing any $e \in {\mathfrak g}^0$
such that $\graded_0(e) = E$, one can show, using induction and
$\cap_j {\mathfrak g}^j=0$, that the maps
$$ X \mapsto  [e, X] - jX \colon {\mathfrak g}^j \rightarrow {\mathfrak g}^{j+1}$$
are well-defined projections, whose kernels constitute a grading of ${\mathfrak g} $ for
which $e$ is a grading element; this in turn gives rise to an isomorphism
$\Phi_e \colon {\mathfrak g} \rightarrow \graded({\mathfrak g})$ which in fact must be a
Lie algebra isomorphism. When we alter the choice of $e$, the isomorphism $\Phi_e$ changes
in a way made precise in \eqref{splittings2} below.

A grading of ${\mathfrak g} $ as we have just described it will be called an {\em
  $E$-splitting} of the filtration. With $E$ and $e$ fixed as above, there is evidently a
one-to-one correspondence between $E$-splittings and the affine space
$e + {\mathfrak g}^1$. The following lemma, proven in \ref{appendix_splittings}, will be
used several times in the sequel.
%
\begin{lemma}
  Let $e \in {\mathfrak g}^0 $ be a lift of some grading element
  $E \in \graded({\mathfrak g})$. Let $G^0$ be any Lie group integrating
  ${\mathfrak g}^0 $, and let $G^1 \subset G^0$ be the connected Lie subgroup with Lie
  algebra ${\mathfrak g}^1$. Then:
  \begin{conditions}
  \item Under the adjoint action, $G^1$ acts freely and transitively on the affine space
    $e + {\mathfrak g}^1\subset {\mathfrak g}^0$, and consequently on $E$-splittings of
    ${\mathfrak g} $. In particular, $G^1$ is a simply connected (nilpotent) Lie group, so
    that the infinitesimal adjoint action of ${\mathfrak g}^1$ on ${\mathfrak g} $ has a
    unique globalization $\Adjoint\colon G^1 \rightarrow \automorphism({\mathfrak g})$
    extending the adjoint action already mentioned.\label{splittings1}.

  \item For any $e_1, e_2 \in e + {\mathfrak g}^1$,
    $\Phi_{e_2} = \Adjoint_g\circ \Phi_{e_1}$, where $g \in G^1$ is the unique element for
    which $e_2 = g\cdot e_1$. \label{splittings2}
  \end{conditions}
\end{lemma}

\subsection{Weyl structures for filtered Lie algebroids}\label{weyl}
A {\df Weyl structure} on a filtered Lie algebroid $A$ will be any section $e$ of $A^0$
such that $E = \graded_0(e)$ is a fibre-wise grading section of $\graded(A^0)$, and such that
\begin{equation}
  [X, e] \in A^1 \enspace\text{for all $X \in A$.}\label{w2}
\end{equation}
If $A$ is firmly filtered, so that we have the action $\adjoint$ of $A$ on $\graded(A^0)$
described in \ref{reduction}, then \eqref{w2} holds if and only if $E$ is
$\adjoint$-invariant.

Note that we do not require that $E$ be a fibre-wise grading section of
$\graded(A)$, only of $\graded(A^0)$ (but see also \ref{compat2}).

\begin{theorem}
  Let $A$ be a filtered Lie algebroid admitting a Weyl structure $e$ and let $B \subset A$
  denote the isotropy of $e$ under the canonical representation of $A$ on its anchor
  kernel, i.e., the kernel of the map
  $X \mapsto [X, e] \colon A \rightarrow A^1$. Then:
  \begin{conditions}
  \item $A$ is firmly filtered, and in particular, $A^1 \subset A$ is an ideal (i.e.,
    $\Gamma(A^1) \subset \Gamma(A)$ is an ideal).\label{pq1}
  \item $A$ is the internal semidirect product of $B$ and $A^1$.\label{pq2}
  \item The grading $A^0=\oplus_{i\ge0}{\mathfrak a}_i$ associated with $e$ is invariant
    under the canonical representation of $B$. In particular, the corresponding
    isomorphism $A^1\cong \graded(A^1)$ interwines the representation of $B$ on $A^1$ with
    the representation of $A/A^1\cong B$ on $\graded(A^1)$, so that $A$ is the external
    semidirect product of $A/A^1$ and $\graded(A^1)$.\label{pq3}
  \end{conditions}
\end{theorem}

\begin{proof}
  The isotropy $B \subset A$ is a subalgebroid, by Jacobi's identity. We first prove
  \begin{equation}
     A = B \oplus A^1, \label{y7}
  \end{equation}
  as vector bundles. It suffices to show that the restriction $r$ of
  $X \rightarrow [X, e]$ to $A^1$ is an isomorphism. However, from the filtration property
  F1, this restriction is filtration-preserving, and the corresponding map
  $\graded(A^1) \rightarrow \graded(A^1)$ sends $\xi \in \graded(A^1)_j$ to $j \xi$,
  because $\graded(e)$ is a fibre-wise grading section of $\graded(A^0)$. As this is
  obviously an isomophism, $r$ must be too.

  Next we prove the grading $A^0=\oplus_{i\ge0}{\mathfrak a}_i$ associated with $e$ is
  invariant under the canonical representation of $B$. To that end, suppose
  $Y \in \Gamma({\mathfrak a}_j)$ for some $j\ge 0$ and let $X \in \Gamma(B)$ be
  arbitrary. The Jacobi identity for the bracket on $A$ gives us
  \begin{equation*}
    [X, [e, Y]] = [[X, e], Y] + [e, [X, Y]].
  \end{equation*}
  On the other hand, applying $\adjoint_X$ to both sides of the
  identity $[e, Y] = jY$, we have
  \begin{equation*}
     [X, [e, Y]] = j\adjoint_XY.
  \end{equation*}
  Combining these two results, we obtain
  \begin{equation*}
    [[X, e], Y] + [e, \adjoint_XY] = j \adjoint_XY.
  \end{equation*}
  But $[X, e] = 0 $, because $X \in \Gamma(B)$. We conclude that
  $\adjoint_XY \in \Gamma({\mathfrak a}_j)$.

  Having proven, in particular, that $B$ preserves the filtration of $A^0$, to establish
  \eqref{pq1} it suffices, in view of \eqref{y7}, to show that the filtration of $A^0$ is
  invariant under the canonical action of itself, which follows immediately from F1.

  Claim \eqref{pq2} follows from \eqref{pq1} and \eqref{y7}. The first part of \eqref{pq3}
  has already been addressed, and the rest now follow easily.
\end{proof}

\subsection{Filtered Lie groupoids and Weyl structures}\label{weyl2}
We will call a Lie groupoid ${\mathcal G} $ over $M$ {\df filtered} if its associated Lie
algebroid $A$ is filtered. Now ${\mathcal G} $ acts on elements of the anchor kernel of
$A$, and so on elements of $A^0$, although this action need not preserve $A^0$ in
general. This action will be denoted $\Adjoint$. The Lie groupoid is {\em firmly filtered}
if $A^j$ is ${\mathcal G}$-invariant under $\Adjoint$, for all $j \ge 0$. The induced
action of ${\mathcal G} $ on $\graded(A^0)$ will also be denoted $\Adjoint$.

We call $e \in \Gamma(A)$ a {\df Weyl structure} for ${\mathcal G} $ if it is a Weyl
structure for $A$, and if for every arrow $g \in {\mathcal G}$ with source $m \in M$ we
have
\begin{equation}
\Adjoint_g e(m) - e(gm) \in A^1(gm),  \label{iden}
\end{equation}
where $gm$ denotes the target of $g$. If ${\mathcal G} $ is firmly filtered, so that
${\mathcal G} $ acts on $\graded(A^0)$, then this condition holds if and only if $E$ is
${\mathcal G}$-invariant under $\Adjoint$, in which case the requirement \eqrefs{weyl}{w2}
is automatic.
% \begin{proposition}
%   If $A^1$ is $\Adjoint$-invariant, then every Weyl structure for ${\mathcal G} $ is a
%   also a Weyl structure for its Lie algebroid $A$.
% \end{proposition}
% \begin{proof}
%   Let $X \in A$, and suppose $t \mapsto g_t \in {\mathcal G} $ is a smooth path of arrows
%   with common source $m \in M$ with $g_0 =m $ and $\dot g_0 =X$. Then we may write
%   \begin{equation*}
%     [X, e](m) = \frac{d}{dt} \Adjoint_{g_t}^{-1}e(g_tm)\,\Big|_{t=0}.
%   \end{equation*}
%   Applying the identity \eqref{iden}, we can write this as
%     \begin{equation*}
%       [X, e](m) = \frac{d}{dt}\Big(\,e(m) + \Adjoint_{g_t}^{-1}X^1(g_tm)\,\Big)\Big|_{t=0} =
%        \frac{d}{dt}\Adjoint_{g_t}^{-1}X^1(g_tm)\,\Big|_{t=0},
%   \end{equation*}
%   where $X^1$ is some section of $A^1$. The $\Adjoint$-invariance of $A^1$ now shows that
%   $[X, e](m) \in A^1(m)$.
% \end{proof}

To navigate a relatively mild obstruction to obtaining semidirect products in the groupoid
case, we will say the filtration of $\mathcal G$ is {\df weakly firm} if $A^1$ is
${\mathcal G} $-invariant under $\Adjoint$. Evidently a firm filtration of
${\mathcal G} $ is weakly firm.
\begin{proposition}
  Suppose ${\mathcal G} $ is filtered and admits a Weyl structure. Then if $\mathcal G$
  has connected source fibres, it is weakly firm.
\end{proposition}
\begin{proof}
  Since any Weyl structure for $\mathcal G$ is a Weyl structure for its Lie algebroid $A$,
  $A^1 \subset A$ must be an ideal, by Theorem \eqrefs{weyl}{pq1}. The result now follows
  from the connectedness assumption.
\end{proof}

We can now state and prove a global version of Theorem \ref{weyl}.

\begin{theorem}
  Let ${\mathcal G} $ be a weakly firm, filtered Lie groupoid with Lie algebroid $A$. Let
  ${\mathcal G}^1 \subset {\mathcal G} $ denote the source-connected Lie subgroupoid with
  Lie algebroid $A^1$ and suppose ${\mathcal G} $ admits a Weyl structure $e$. Then, if
  ${\mathcal H} \subset {\mathcal G} $ denotes the isotropy subgroupoid of $e$ under the
  canonical representation of ${\mathcal G} $:
  \begin{conditions}
  \item ${\mathcal G}^1 \subset {\mathcal G} $ is normal.\label{pr1}
  \item ${\mathcal G} $ is the internal semidirect product of ${\mathcal H} $ and
    ${\mathcal G}^1$.\label{pr2}
  \item The grading $A^0=\oplus_{i\ge0}{\mathfrak a}_i$ associated with $e$ is invariant
    under the canonical representation of ${\mathcal H}$. \label{pr3}
  \item ${\mathcal G} $ is firmly filtered.\label{pr4}
  \end{conditions}
\end{theorem}
\begin{proof}
  Since we assume ${\mathcal G} $ is weakly firm, ${\mathcal G}^1 \subset {\mathcal G} $
  must be normal, because ${\mathcal G}^1$ is source-connected, by definition. This proves
  \eqref{pr1}.

  To prove \eqref{pr2} we now show
  each arrow $g \in {\mathcal G}$, beginning at some $m \in M$, can be written uniquely as
  $g = kh$, with $k \in {\mathcal G}^1$ and $h \in {\mathcal G} $. Indeed, from
  \eqref{iden} we have $\Adjoint_ge(m) \in e(gm) + A^1(gm)$. But, according to Lemma
  \ref{splittings}, the Lie group ${\mathcal G}^1_{gm}$ (the source=target fibre of
  ${\mathcal G}^1$ at $gm \in M$) acts freely and transitively on the affine space
  $e(gm) + A^1(gm)$. Therefore there exists a unique $k \in {\mathcal G}^1$ such that
  $\Adjoint_k e(gm) = \Adjoint_ge(m)$. But this means
  $\Adjoint_{k^{-1}g}e(m) = e(gm) = e(k^{-1}gm)$, i.e., $k^{-1}g \in {\mathcal H} $.

  To prove \eqref{pr3}, suppose $Y \in \Gamma({\mathfrak a}_j)$ for some $j\ge 0$ and let
  $h \in {\mathcal H} $ be an arbitrary arrow beginning at some $m \in M$. Since the
  bracket on $A^0$ is ${\mathcal G}$-invariant under $\Adjoint$, the grading identity
  $[e, Y](m) = jY(m)$ gives us $[\Adjoint_he(m), \Adjoint_hY(m)] = j\Adjoint_hY(m)$. From
  the definition of ${\mathcal H} $, we conclude that
  $[e(hm), \Adjoint_hY(m)] = j\Adjoint_hY(m)]$, establishing that
  $\Adjoint_hY(m) \in {\mathfrak a}_j$, as required.

  In view of the fact that ${\mathcal G} = {\mathcal G}^1 {\mathcal H} $, and having
  established \eqref{pr3}, it suffices, in proving \eqref{pr4}, to show that the
  filtration of $A^0$ is invariant under the action of ${\mathcal G}^1$. But $A^0$ is
  already invariant under the action of $A^1$ by the filtration property F1. Since
  ${\mathcal G}^1$ is source-connected, the filtration of $A^0$ is therefore invariant
  under the action of ${\mathcal G}^1$.
\end{proof}

\subsection{The existence of Weyl structures}\label{existence_weyl}
We can now characterize those filtered Lie algebroids/groupoids admitting Weyl
structures.

\begin{theorem}\mbox{}
  \begin{conditions}
  \item A filtered Lie algebroid $A$ admits a Weyl structure if and only if it is firmly
    filtered and $\graded(A^0)$ admits an $\adjoint$-invariant fibre-wise grading section $E$.
    \label{ew1}
\item A weakly firm, filtered Lie groupoid ${\mathcal G} $ admits a Weyl
  structure if and only if it is firmly filtered and $\graded(A^0)$ admits an
  $\Adjoint$-invariant fibre-wise grading section $E$.\label{ew2}
  \end{conditions}
\end{theorem}
\begin{proof}
  Given a firmly filtered Lie algebroid $A$ with $\adjoint$-invariant fibre-wise grading
  section $E$, any lift of $E$ to a section $e \in \Gamma(A^0)$ is readily seen to be a
  Weyl structure. The converse claim follows from Theorem \ref{weyl}. The groupoid case is
  similar, with the converse part following from Theorem \ref{weyl2}.
\end{proof}

\section{Cartan connections}\label{cartan}

This section recalls the notion of a Cartan connection on a Lie groupoid, or a Lie
algebroid, and summarizes properties needed in the sequel.

\subsection{Cartan connections on a Lie groupoiod} \label{cartan} A {\df Cartan
  connection} on a Lie groupoid ${\mathcal G} $ over $M$ is an $n$-plane field
$D \subset T {\mathcal G}$ arising from some right inverse
${\mathcal G} \rightarrow J^1 {\mathcal G} $ for the canonical projection
$J^1 {\mathcal G} \rightarrow {\mathcal G} $, which is necessarily uniquely determined by
$D$ \cite[Appendix A]{Blaom_12}, \cite{Blaom_16b}. Here $J^1 {\mathcal G} $ denotes the Lie
groupoid of one-jets of bisections of ${\mathcal G} $. Of course, we are understanding the
right inverse to be a Lie groupoid morphism. This is equivalent to requiring that that
$D \subset TM$ be multiplicatively closed, in the sense of
\cite{Salazar_13,Crainic_etal_15}, but this fact will not play a role here.

The usual adjoint action of a Lie group on its Lie algebra is generalized by an action of
$J^1 {\mathcal G} $ on the Lie algebroid $A$ of ${\mathcal G} $, denoted
$\Adjoint$. Composing with the right inverse ${\mathcal G} \rightarrow J^1 {\mathcal G} $,
we obtain a representation $\bar D$ of ${\mathcal G}$ on $A$. This representation extends
the canonical one of ${\mathcal G} $ on elements of the anchor kernel of $A$ (which we
have been also denoting by $\Adjoint$ in Section \ref{semidirect}). In the transitive
case, every such representation is the representation $\bar D$ corresponding to some {\df
  unique} Cartan connection $D$.

The prototype of a Cartan connection is the horizontal $n$-plane field $D$ on
an action groupoid $G \times M$, for a smooth action of some Lie group $ G $ on
$M$, which need not be transitive. In this case

\begin{equation}
  \bar D_{(g, m)}(\xi, m) = (\Adjoint_g \xi, g\cdot m).\label{adjoint}
\end{equation}

Cartan connections on Lie groupoids generalize to the intransitive case the classical
Cartan connections introduced by Ehresmann \cite{Ehresmann_51} and in widespread use, when
the latter notion is ``stripped of model data'' \cite{Blaom_16b}.

\subsection{Infinitesimal Cartan connections}\lab{hsd}%
Infinitesimalizing a Cartan connection, viewed as a Lie groupoid
morphism ${\mathcal G} \rightarrow J^1 {\mathcal G}$, we obtain a splitting $s
\colon A \rightarrow J^1 A $ of the exact
sequence
\begin{equation*}
  0\longrightarrow T^*\!M\otimes A\longrightarrow J^1A\longrightarrow A\longrightarrow 0
\end{equation*}
in the category of Lie algebroids. Here $J^1 A $ denotes the Lie algebroid of one-jets of
sections of the Lie algebroid $A $ of ${\mathcal G}$. In the category of vector bundles, splittings
$s$ of the above sequence are in one-to-one correspondence with Koszul connections
$\nabla$ on $A$; this correspondence is given by\footnote{According to our sign
  conventions, the inclusion $T^*\!M \otimes A \rightarrow J^1 A $ induces the following
  map on sections: $df \otimes X \mapsto f J^1 X - J^1 (f X)$.}
\begin{equation}
  sX = J^1 X + \nabla X; \qquad  X \in \Gamma(A). \mathlab{cspan}
\end{equation}
When $s \colon A \rightarrow J^1 A $ is a morphism of Lie algebroids, as is the case here,
then $\nabla $ is called a {\df Cartan connection} on $A$ \cite{Blaom_06}, or an {\df
  infinitesimal Cartan Connection} when we want to distinguish it from the corresponding
global notion. Evidently, the notion makes sense for arbitrary Lie algebroids, not just
those integrating Lie groupoids. According to \cite{Blaom_06}, $\nabla $ is a Cartan
connection on ${\mathfrak g} $ if and only if
\begin{equation*}
  \nabla_V[X, Y] = [\nabla_VX, 
        Y]+[X,\nabla_VY]-\nabla_{\bar\nabla_XV}Y+\nabla_{\bar\nabla_YV}X,
\end{equation*}
for all vector fields $V$ on $M$ and sections $X, Y$ of $A$. Here $\bar\nabla$ denotes the
associated $A$-connection on $TM$ defined by
\begin{equation}
 \bar\nabla_XV= \#\nabla_VX +[\#X, V],
\end{equation}
an $A$-representation when $\bar\nabla $ is Cartan. Here $\# \colon A \rightarrow TM$
denotes anchor. Often the following equivalent condition is more useful:
\begin{equation}
  (\nabla_V\torsion \bar\nabla)(X,Y) =
 \curvature \nabla (\#X, V)Y -  \curvature \nabla (\#Y, V) X.
\end{equation}

Given a Cartan connection $\nabla $ on some Lie algebroid $A$, we can compose the
corresponding splitting $s \colon A \rightarrow J^1 A$ with the adjoint action of $J^1 A$
on $ A$, $\adjoint_{J^1X}Y = [X, Y]$, to obtain a self-representation $\bar \nabla $ of
$A$. Explicitly, we have

\begin{equation}
 \bar\nabla_XY = \nabla_{\#Y}X +[X, Y].\label{zero}
\end{equation}
The anchor $\# \colon A \rightarrow TM$ becomes $\bar\nabla$-equivariant.

The representation $\bar\nabla$ of $A$ on iteself extends the canonical representation of
$A$ on elements of its anchor kernel (denoted $\adjoint$ in Section \ref{semidirect}). In
the transitive case, any such representation is $\bar\nabla $ for some uniquely determined
Cartan connection $\nabla $.

\subsection{Curvature as the obstruction to symmetry}\label{curvature}
The {\df curvature} of $\nabla $ is defined as usual. This notion of curvature is not to
be confused with the notion of curvature for classical Cartan connections, as this is
generally understood, which relies on the specification of model data. We will call this
latter notion {\df model curvature}. The Lie algebroid analogue of model data is
introduced in Section \ref{modeldata}; model curvature is defined in
\ref{model_curvature}. For further context, see \cite{Blaom_16b}. The {\df torsion} is
defined by
\begin{equation*}
  \torsion \bar\nabla (X, Y) = \bar\nabla_XY - \bar\nabla_YX - [X, Y] =
  \nabla_{\#Y}X - \nabla_{\#X}Y + [X, Y].
\end{equation*}

The prototype of an infinitesimal Cartan connection is the trivial flat connection
$\nabla $ on an action algebroid ${\mathfrak g} \times M$, for some infinitesimal action
of a Lie algebra $\mathfrak g$ on $M$. If we identify ${\mathfrak g} $ with the constant
($\nabla$-parallel) sections of ${\mathfrak g} \times M$, then its bracket is
$\torsion\bar\nabla(\,\cdot\,,\,\cdot\,)$. Indeed, let
$\tau \in \Gamma(\wedge^2 A^* \otimes A)$ denote the trivial extension of the bracket on
${\mathfrak g} $ to a fibre-wise bracket on $A = {\mathfrak g} \times M$, and recall that
the Lie algebroid bracket on ${\mathfrak g} \times M$ is given by
\begin{equation}
  [X,Y]=\nabla_{\#X}Y-\nabla_{\#Y}X +\tau(X,Y). \label{one}
\end{equation}
Then we can write
\begin{equation}
    [X,Y]=\bar\nabla_XY-\bar\nabla_YX - \tau(X,Y),\label{two}
\end{equation}
from which $\tau = \torsion \bar\nabla $ immediately follows.

\begin{theorem}[\cite{Blaom_06}]
  If $M$ is simply connected, then every Lie algebroid $A$ equipped with a flat Cartan
  connection $\nabla$ is naturally isomorphic to an action algebroid
  ${\mathfrak g} \times M$, where ${\mathfrak g}\subset \Gamma(TM) $ is the Lie algebra of
  $\nabla $-parallel sections, equipped with bracket $\torsion \bar\nabla$.
\end{theorem}
\noindent More informally, we have:
\begin{aside}
  A Lie algebroid equipped with a Cartan connection is an infinitesimal symmetry deformed
  by curvature; the torsion of the connection is the deformed symmetry Lie bracket.
\end{aside}

If $\curvature \nabla =0$ but $M $ is not simply connected, then $A$ is a Lie algebra
action with monodromy twist \cite{Blaom_16b}. Lie groupoids with Cartan connnections,
whose infinitesimalizations are flat, define pseudogroups of symmetries of $M$, and under
an additional completeness assumption, twisted Lie group actions \cite{Blaom_16,
  Blaom_16b}. In the transitive case, one can recover a locally homogeneous structure,
assuming ``geometric closure'', if not global homogeneity up to cover \cite{Blaom_13}.

\subsection{Cartan connections compatible with a filtration}\label{compat}
An infinitesimal Cartan connection $\nabla$ on a filtered Lie algebroid $A$ will be {\df
  compatible} with the filtration if the filtration is invariant under the associated
representation $\bar\nabla $. Similarly, a Cartan connection $D$ on a filtered Lie
groupoid ${\mathcal G} $ with Lie algebroid $A$ will be said to be {\df compatible} with
the filtration provided the filtration of $A$ is invariant under the associated
representation $\bar D$, which in turn implies compatability of the corresponding
infinitesimal Cartan connection $\bar\nabla$.

\begin{proposition}
  If a filtered Lie groupoid (Lie algebroid) admits a compatible (infinitesimal) Cartan
  connection, then it is firmly filtered.
\end{proposition}
\begin{proof}
  Suppose $\nabla $ is a Cartan connection on a filtered Lie algebroid $A$. Then
  $\bar\nabla $ extends the canonical representation $\adjoint$ of $A$ on its anchor
  kernel, which contains $A^0$. So $\adjoint_X Y \equiv [X, Y] = \bar\nabla_XY$ for
  $X \in A$ and $Y \in A^0$. If $\nabla $ is compatible with the filtration, then $A^0$
  and the filtration of $A^0$ must therefore be $\adjoint$-invariant. The groupoid case is
  similar.
\end{proof}

\subsection{Filtered Klein geometries}\label{homogeneous}
We now prove some claims made in \ref{eg2} regarding the homogeneous space $M=G/P$
described there. Let ${\mathcal G} = G \times M$ be the action groupoid, equipped with its
canonical Cartan connection $D$, so that ${\mathcal G}$ acts on
$A = {\mathfrak g} \times M$ according to \eqrefs{cartan}{adjoint}. Since the filtration
of ${\mathfrak g} $ is $P$-invariant, the filtration of $A$ is $\bar D$-invariant.

Since $A$ is transitive and $A^0$ is the anchor kernel, we have, as vector bundles,
\begin{equation*}
  \graded(A) \cong \graded(TM)\oplus \graded(A^0).
\end{equation*}
As in \ref{curvature}, we denote the fibre-wise Lie bracket on $A={\mathfrak g} \times M$
coming from ${\mathfrak g} $ by $\tau $, which coincides with the torsion of the
infinitesimalization $\nabla $ of the $D$ (the canonical flat Koszul connection on the
trivial bundle ${\mathfrak g} \times M$). We note in passing that $\tau$ is evidently
$\bar D$-invariant, and so $\bar\nabla$-invariant. When we use $\tau$ to view
$\graded(A)$, $\graded(A^0)$ or $\graded(TM)$ as Lie algebra bundles, we shall denote them
by $\graded_\tau(A)$, $\graded_\tau(A^0)$ and $\graded_\tau(TM)$.

\begin{proposition}
  The vector bundle filtration of $A = {\mathfrak g} \times M$ defined in \ref{eg2} is a
  firm Lie algebroid filtration, so that ${\mathcal G} = G \times M$ is a firmly filtered
  Lie groupoid. Furthermore, equipping $\graded(A)$ with the corresponding Lie algebra
  bundle structure (Proposition \ref{filtered}), we have
  \begin{equation*}
    \label{decomposition}
    \graded(A) = -\graded_\tau(TM)\oplus \graded_\tau(A^0)\quad
    \text{(Lie algebra bundle direct sum)}
  \end{equation*}
  where -$\graded_\tau(TM)$ is $\graded_\tau(TM)$ equipped with
  negative bracket.
\end{proposition}

\begin{proof}
  Axiom F2 holds by construction. Let $X\in \Gamma(A^j)$ and $Y\in \Gamma(A^k)$ be
  given. Then we have:
  \begin{align}
  \bar\nabla_X Y &= [X, Y] \quad\text{if $k \ge 0$}\label{gamma}\\
  \bar\nabla_X Y &= \tau(X, Y) \quad\text{if $j \ge 0$}\label{delta}\\
  [X, Y] &= \tau(X, Y) \quad\text{if $j \ge 0$ and $k \ge 0$} \label{alpha}\\
  [X, Y] &= -\tau(X, Y) \mod A^{j + k + 1} \quad\text{if $j < 0$ and $k < 0$}\label{beta}\\
  [X, Y] &\in \Gamma(A^{j + k + 1}) \quad\text{if $j < 0$ and $k \ge 0$}.\label{blue}
  \end{align}
  Equation \eqref{gamma} follows from \eqrefs{hsd}{zero} and F2. If $j\ge 0$, then
  Equation \eqrefs{curvature}{one} reads $[X, Y] = - \nabla_{\#Y}X +\tau(X, Y)$, because
  $\#A^j=0$ in that case. This and \eqrefs{hsd}{zero} delivers \eqref{delta}.  Claim
  \eqref{alpha} follows immediately from the preceding two.

  If $j < 0$, then $j + k + 1 \le k$, and so
  $A^k \subset A^{j + k+ 1}$. Similary, if $k < 0$, then
  $A^j \subset A^{j + k + 1}$. Since the filtration is
  $\bar\nabla$-invariant, \eqref{beta} now follows from \eqrefs{curvature}{two}.

  As already mentioned, if $j <0 $, then $A^k \subset A^{j+k+1}$. On
  the other hand, if additionally $k \ge 0$, then
  $[X, Y] = \bar\nabla_XY$, by \eqref{gamma}. So \eqref{blue} is a
  consequence of the $\bar\nabla $-invariance of $A^k$.

  To establish axiom F1 for filtered Lie algebroids, we need to show
  $[X, Y] \in \Gamma(A^{j+k})$. Equations \eqref{alpha} takes care of the case $j\ge 0$
  and $k \ge 0$. Equation \eqref{beta} takes care of the case $j < 0$ and $k < 0$ because
  $A^{j+k+1} \subset A^{j+k}$ is always true. The remaining case is $j < 0$ and $k \ge 0$
  (and the case $j \ge 0$ and $k < 0$, which will follow by symmetry). But as
  $A^{j + k + 1} \subset A^{j + k}$, this already follows from \eqref{blue}.

  The second statement in of the theorem follows easily from \eqref{alpha}, \eqref{beta},
  and \eqref{blue}.

  Since the filtration is $A$ is $\bar D$-invariant, $D$ is compatible with the
  filtration, and so firmness follows from Proposition \ref{compat}.
\end{proof}

\section{Model data}\label{modeldata}
Let $A$ be a filtered Lie algebroid $A$. Since $A^0$ lies in the anchor kernel, it is a
Lie algebra bundle. We we call an extension $\tau$ of this Lie algebra bundle structure,
to one on $A$, {\df model data}.

\subsection{Compatibility}\label{compat2}
Model data $\tau$ is {\df compatible} with a filtration of $A$
if this is also a filtration with respect to $\tau$, i.e., $\tau(X, Y) \in A^{j+k}$
whenever, $X \in A^j$ and $Y \in A^k$ have a common base point. We say $\tau $ is
compatible with an infinitesimal Cartan connection $\nabla $ (Cartan connection $D$) if
$\tau $ is $\bar \nabla $-invariant (respectively $\bar D$-invariant). If, additionally,
\begin{equation}
  \bar\nabla_XY = \tau(X, Y)\quad \text{for all $X \in A^0$, $Y \in \Gamma(A),$}\mathlab{strong}
\end{equation}
then we say $\nabla$ is {\em strongly compatible} with $\tau$. We say a Cartan connection
$D$ is {\em strongly compatible} with $\tau$, if its infinitesmialization $\nabla $ is
strongly compatible with $\tau $. Theorem \ref{parabolic} stated later shows that
compatibility already implies strong compatibility, for model data relevant to parabolic
geometry.\footnote{Strong compatibility is built in to the definition of classical Cartan
  connections with model data, but we won't always need it here.} We say that $\tau$ is
{\df compatible} with a Weyl structure $e \in \Gamma(A^0)$ if $E = \graded_0(e)$ is a
fibre-wise grading element of $\graded_\tau(A)$ (and not just of
$\graded(A^0)=\graded_\tau(A^0)$).

If $A$, ${\mathfrak g}$, and $\tau$ have the meanings given for filtered Klein geometries
in \ref{homogeneous}, then $\tau$ is model data compatible with the filtration of $A$
induced by that of ${\mathfrak g} $, and the canonical flat connection $\nabla $ on
$A = {\mathfrak g} \times M$ is strongly compatible with $\tau $, by
\eqrefs{homogeneous}{delta}.

Some of our notions of compatability are summarized in Table 1.
\begin{table}[]
\centering
  \begin{tabular}{@{}|c|c|c|c|@{}}
  \toprule
  compatibility of $\downarrow$ with $\rightarrow$ & $\tau$   & $\nabla$                        &  $e$\\ \midrule
  filtration & $\tau(A^j, A^k) \subset A^{j+k}$ &  $A^j$ is $\bar\nabla$-invariant $\forall j$ & -   \\
  $\tau$     &                                 &  $\tau $ is $\bar\nabla $-invariant  & $\graded_0(e)$ grades $\graded_\tau(A)$ \\
  $\nabla $  &                                 &                                      &  $e$ is $\bar\nabla $-invariant \\  \bottomrule
\end{tabular}
\vspace{0.75\baselineskip}
\caption{Summary of compatability conditions for structures on a filtered Lie algebroid
  $A$. Here $\tau $ model data, $\nabla $ an infinitesimal Cartan connection, and $e$ a
  Weyl structure.}
\end{table}

\subsection{Regularity}\label{regularity}
In the special transitive case, where $\oplus_{i<0}\graded(A)_i \cong \graded(TM)$, model
data $\tau $ compatible with the filtration of $A$ will be said to be {\df regular} if the
induced bracket on $\graded(TM)$ is minus one times the canonical bracket (the one defined
by the Lie algebroid bracket): $\graded_\tau(TM) = -\graded(TM)$. According to Proposition
\ref{homogeneous}, the canonical model data $\tau $ for a filtered Klein geometry is
always regular.

\subsection{Gradings associated with model data}\label{model_data_gradings}
Let $A$ be a filtered Lie algebroid admitting a Weyl structure $e$. Then, according to
Theorem \ref{weyl}, we have an associated grading $A^0=\oplus_{i\ge 0}{\mathfrak a}_i$
that is invariant under the action of the isotropy subalgebroid $B \subset A$ of $e$. If
$e$ is compatible with model data $\tau$ compatible with the filtration of $A$, then this
grading extends to a grading $A = \oplus_i {\mathfrak a}_i $, where
$X \in {\mathfrak a}_j$ if and only if $\tau(e, X) = jX$.

\begin{proposition} In the scenario above, let $\nabla $ be any infinitesimal Cartan
  connection compatible with $\tau $. Then:
  \begin{conditions}
  \item The grading $A = \oplus_i {\mathfrak a}_i$ is invariant under the representation
    $\bar\nabla $ of $B$ on $A$.    \label{u4}
  \item For any $X \in \Gamma({\mathfrak a}_k)$, $k\ge 0$, and
    $Y \in \Gamma({\mathfrak a}_j)$, $j \in {\mathbb Z} $, one has
    $\bar\nabla_XY \in \Gamma({\mathfrak a}_j \oplus {\mathfrak a}_{j+k})$.\label{u3}
%  \item $\nabla $ is compatible with the filtration of $A$. \label{u7}
  \end{conditions}
  If ${\mathcal G} $ is a weakly firm, filtered Lie groupoid with Lie algebroid $A$, $e$
  is a Weyl structure for ${\mathcal G} $ compatible with the model data $\tau $, and $D$
  is a Cartan connection on ${\mathcal G} $ compatible with $\tau $, then:
  \begin{conditions}
  \item The grading $A=\oplus_i {\mathfrak a}_i$ associated with $e$ is invariant under
    the representation $\bar D $ of the isotropy subgroupoid
    ${\mathcal H}\subset {\mathcal G} $ of $e$.\label{uu4}
%  \item $D$ is compatible with the filtration of $A$.\label{uu7}
  \end{conditions}
\end{proposition}

The reader is warned that $B$ does not generally align itself with the complement
$\oplus_{i\le 0}{\mathfrak a}_i$ of $A^1=\oplus_{j\ge 1}{\mathfrak a}_i$ in $A$ defined by
the grading.  We may only conclude that $B$ is the graph of some linear function
$\oplus_{j\le 0} {\mathfrak a}_i \rightarrow A^1$ that vanishes on ${\mathfrak a}_0$,
which is necessarily contained in $B$.

\begin{proof}
  Our proofs of \eqref{u4} and \eqref{u3} are variations on an argument appearing in the
  proof of Theorem \ref{weyl}.  Indeed, suppose that $Y \in \Gamma({\mathfrak a}_j)$ for
  some $j$ and let $X \in \Gamma(B)$ be arbitrary. Then, applying $\bar\nabla_X$ to both
  sides of the identity $\tau(e, Y) = jY$, and appealing to the $\bar\nabla$-invariance of
  $\tau$, we have
  \begin{equation*}
    \tau(\bar\nabla_{X}e, Y) + \tau(e, \bar\nabla_XY) = j \bar\nabla_XY.
  \end{equation*}
  But $\bar\nabla_{X} e = [X, e]$, because $e$ is in the anchor kernel, while
  $[X, e] = 0 $, because $X \in \Gamma(B)$. We conclude that
  $\bar\nabla_XY \in \Gamma({\mathfrak a}_j)$. This proves \eqref{u4}.

  Turning to \eqref{u3}, write $\bar\nabla_XY=\sum_iW_i$ for some
  $W_i \in \Gamma({\mathfrak a}_i)$, which implies
  \begin{equation*}
    \tau(e, \bar\nabla_XY) = \sum_iiW_i.
  \end{equation*}
  On the other hand, $\tau(e, Y) = jY$, which, after applying $\bar\nabla_X$ to both
  sides, and appealing to the $\bar\nabla$-invariance of $\tau $, becomes
  \begin{equation*}
    \tau(\bar\nabla_X e, Y) + \tau(e, \bar\nabla_XY) = j\sum_iW_i.
  \end{equation*}
  As $X$ and $e$ are sections of $A^0$, we have
  $\bar\nabla_X E = [X, e] = \tau(X, e) = -kX$. From this and the two preceding equations we
  deduce
  \begin{equation*}
    \sum_i(i-j)W_i + k \tau(Y, X) =0.
  \end{equation*}
  Since $A = \oplus_i {\mathfrak a}_i$ is a grading with respect to $\tau$, the last term
  on the left-hand side is a section of $\Gamma({\mathfrak a}_{j+k})$. Projecting onto
  each component ${\mathfrak a}_i$ of the grading $A = \oplus_i {\mathfrak a}_i$, we
  obtain $W_i = 0$ for all $i$, with the possible exceptions $i=j, j+k$, which shows
  $\bar\nabla_XY \in \Gamma({\mathfrak a}_j \oplus {\mathfrak a}_{j+k})$.

  The proof of \eqref{uu4} is similar to that of \eqref{u4}, following from the
  $\bar D$-invariance of $\tau $, the definition of the spaces ${\mathfrak a}_j$, and the
  fact that ${\mathcal H} $ fixes $e$.
\end{proof}

\begin{corollary}
  If a (weakly firm) filtered Lie algebroid (Lie groupoid) admits a Weyl structure
  compatible with some model data $\tau $, then any Cartan connection compatible with
  $\tau $ is compatible with the filtration.
\end{corollary}

\begin{proof}
  If $\nabla $ is the Cartan connection in the Lie algebroid case, then
  $\bar\nabla$-invariance of the filtration of $A$ follows from \eqref{u4} and \eqref{u3},
  and the semi-direct product structure of $A$ (Theorem \eqrefs{weyl}{pq2}).

  For the Lie groupoid case and a global Cartan connection $D$, we know the filtration is
  $\bar\nabla $-invariant, where $\nabla $ is the infinitesimalization of $D$, from the
  Lie algebroid case just proven. In particular, continuing in the notation of the
  proposition, the restriction of $\bar\nabla $ to a representation of $A^1$ on $A$ is
  filtration-preserving. Therefore, the restriction of $\bar D$ to a representation of
  ${\mathcal G}^1$ is filtration-preserving, because ${\mathcal G}^1$ is source-connected.
  Full ${\bar D}$-invariance of the filtration follows from this fact, \eqref{uu4}, and
  the semidirect product structure of $\mathcal G$ (Theorem \eqrefs{weyl2}{pr2}).
\end{proof}

\section{Reconstructing Cartan connections with model data}

\subsection{Reduced geometric data}\label{reduced_geometric}
Let $A$ be a filtered Lie algebroid, $\tau $ filtration-compatible model data, and
$\nabla $ a Cartan connection compatible with $\tau $ and the filtration. Then we have the
filtration-preserving representation $\bar\nabla $ of $A$ on itself, which in turn leads
to a grading-preserving representation of $A$ on $\graded(A)$, a representation that must
preserve the Lie algebra bundle bracket $T = \graded_0(\tau)$. We denote this
representation by $R^\nabla$.

Supposing that $\nabla $ is strongly compatiblity with $\tau$, $A^1$ is here acting
trivially on $\graded(A)$ (because $\tau$ is compatible with the filtration) so we even
get a representation of the Levi reduction $A/A^1$ on $\graded(A)$, also denoted
$R^\nabla$.

If additionally, $A$ admits a Weyl structure $e$ compatible with $\tau $ (in which case,
we can drop the requirement that $\bar\nabla $ be filtration-compatible, as this will
already follow from Corollary \ref{model_data_gradings}) then, by the definition of Weyl
structures, $E = \graded(e)$ will be $R^\nabla$-invariant, because $R^\nabla $ acting on
$\graded(A^0)$ is the canonical representation.

Similarly, suppose that ${\mathcal G} $ is a weakly firmly filtered Lie groupoid with Lie
algebroid $A$, with filtration-compatible model data $\tau $, and a Cartan connection $D$
compatible with $\tau $ and the filtration of $A$. Then, by similar reasoning, we obtain a
representation of ${\mathcal G}$ on $\graded(A)$ that we denote $R^D$, and this
representation preserves the bracket $T = \graded_0(\tau)$. Supposing that $D$ is strongly
compatible with $\tau$, the source-connected Lie subgroupoid
${\mathcal G}^1 \subset {\mathcal G} $ with Lie algebroid $A^1$ is here acting trivially
on $\graded(A)$, giving us a representation of ${\mathcal G}/{\mathcal G}^1$ on
$\graded(A)$, also denoted by $R^D$. If $e $ is a Weyl structure for ${\mathcal G} $
compatible with $\tau$, then $E = \graded_0(\tau)$ will be $R^D$-invariant.

The preceding observations motivate the following:

\begin{definition}
  {\df Reduced geometric data} for a firmly filtered Lie algebroid $A$ is a triple
  $(R, T, E)$ of structures on $\graded(A)$, such that:
  \begin{conditions}
    \item $R$ is a representation of $A/A^1$ on $\graded(A)$ extending the canonical
      representation of $A/A^1$ on $\graded(A^0)$.\label{rgd1}
    \item $T$ is an $R$-invariant, grading-compatible, Lie algebra bundle bracket on
      $\graded(A)$ extending the canonical one on $\graded(A^0)$ (but typically different
      from the canonical bracket on $\graded(A)$).\label{rgd2}
    \item $E$ is an $R $-invariant fibre-wise grading section for the bracket $T$.\label{rgd3}
      \end{conditions}
    {\df Reduced geometric data} for a firmly filtered Lie groupoid ${\mathcal G} $ with
    Lie algebroid $A$ is a triple $(R, T, E)$ of structures on $\graded(A)$, satisfying the
    globalizations of \eqref{rgd1}--\eqref{rgd3}:
    \begin{conditions}
    \item $R$ is a representation of ${\mathcal G}/{\mathcal G}^1$ on $A$ extending the
      canonical representation of $A/A^1$ on $\graded(A^0)$. \label{rgd5}
    \item $T$ is an $R$-invariant, grading-compatible, Lie algebra bundle bracket on
      $\graded(A)$ extending the canonical one on $\graded(A^0)$.
    \item $E$ is an $R$-invariant fibre-wise grading section for the bracket $T$.
  \end{conditions}
\end{definition}

Evidently, reduced geometric data $(R, T, E)$ for a firmly filtered Lie groupoid
infinitesimalizes to geometric data $(dR, T, E)$ for the associated Lie algebroid. In the
special transitive case, where, we say $T$ is {\df regular} if the restriction of $T$ to
$\oplus_{i<0}\graded(A)_i \cong \graded(TM)$ is minus one times the canonical bracket on
$\graded(TM)$.

\begin{example}[Projective structures]
  Define $A \subset J^2(TM)$ and $T$ as in \ref{projective}. Then $A/A^1 \cong J^1 (TM)$
  acts on $TM$ by adjoint action, and so on
  $\graded(A) = TM \oplus \gl(TM) \oplus T^*\!M$. Choosing $R $ to be this
  representation, we get reduced geometric data $(R, T, E)$ if we choose
  $E(m) = -\identity$ (the {\em unique} fibre-wise grading section for $T$).

  Let ${\mathcal G} \subset J^2(M \times M)$ be the Lie groupoid of 2-jets $J^2_m\phi$ of
  local diffeomorphisms $\phi $ preserving the projective structure to second order
    at $m$. Then ${\mathcal G} $ has Lie algebroid $A$,
  ${\mathcal G}/{\mathcal G}^1 \cong J^1(M \times M)$ acts on $TM$ by adjoint action, and
  whence on $\graded(A)$. Taking $R$ to be this representation, we obtain reduced
  geometric data $(R, T, E) $ for ${\mathcal G} $.
\end{example}

\begin{proposition}
  Let $(R, T, E)$ be reduced geometric data for some firmly filtered
  Lie algebroid $A$ (or Lie groupoid ${\mathcal G}$). Then any lift $e$ of $E$ through
  $A^0 \rightarrow A^0/A^1$ is a Weyl structure for $A$ (resp. ${\mathcal G}$).
\end{proposition}
\begin{proof}
  In the Lie algebroid case, we see that $E$ must be $\adjoint$-invariant, because it is
  $R$-invariant, and $R$ extends the canonical action of $A/A^1$ on $\graded(A)^0$. So $e$
must be a Weyl structure for $A$. The groupoid case is similar.
\end{proof}

\subsection{Reconstruction in the Lie algebroid case}\label{reconstruction_algebroid}
Addressing the problem of constructing Cartan connections and compatible model data from
reduced geometric data, we have the following important result.
\begin{theorem}
  Let $A$ be a transitive, firmly filtered Lie algebroid whose anchor kernel is $A^0$,
  with reduced geometric data $(R, T, E)$. Then there exists filtration-compatible model
  data $\tau $, a filtration compatible Cartan connection $\nabla $, and a Weyl structure
  $e$, such tha $\tau $, $\bar\nabla$ and $e$ are mututally compatible, $\nabla $ is
  strongly compatible with $\tau $, $R^\nabla=R$, $\graded_0(\tau)=T$, and
  $\graded_0(e) = E$. The Lie algebra bundle structures $\tau $ and $T$ are modeled on the
  same Lie algebra. If $T$ is regular, then so is $\tau $.
\end{theorem}

\begin{proof}
  Let $A$ be a firmly filtered Lie algebroid with reduced geometric data $(R, T, E)$. Let
  $e$ be any section of $A^0$ lifting $E$ through $A^0 \rightarrow A^0/A^1$, so that $e$
  is a Weyl structure, by Proposition \ref{reduced_geometric}. Let $B \subset A$ be
  the kernel of $X \mapsto [X, e]$, which maps isomorphically onto $A/A^1$ under
  $A \rightarrow A/A^1$, by Proposition \eqrefs{weyl}{pq2}. Arbitrarily extend the grading
  $A^0 = \oplus_{i\ge 0} {\mathfrak a}_i$ induced by $e$ to a splitting
  $A = \oplus_i {\mathfrak a}_i$ of $A$, as a filtered vector bundle, and let
  $\Phi \colon A \rightarrow \graded(A)$ denote the induced isomorphism. This isomophism
  must intertwine some some representation $\rho^B $ of $B$ on $A$ with the representation
  $R$ of $A/A^1$ on $\graded(A)$:
  \begin{equation}
    \Phi(\rho^B_XZ) = R_{X\modulo A^1}\,\Phi(Z);
    \qquad X \in \Gamma(B),\, Z\in \Gamma(A).\label{intertwining}
  \end{equation}
  To prove the theorem, we establish the following:
  \begin{conditions}
  \item $\phi $ pulls $T$ back to
    filtration-compatible model data $\tau $ on $A$.\label{re1}
  \item Assuming $A$ is transitive and $A^0$ is the anchor kernel, there exists a unique
    Cartan connection $\nabla $ on $A=B \oplus A^1$ such that
    \begin{equation*}
      \bar\nabla_{X + Y}Z = \rho^B_XZ + \tau(Y,Z);
      \qquad X \in \Gamma(B),\,Y \in \Gamma(A^1),\, Z \in \Gamma(A).
    \end{equation*}\label{re2}
  \item In this case, $\nabla $, $\tau $ and $e$ are mutually compatible, and $\nabla $ is
    strongly compatible with $\tau$.\label{re3}
  \end{conditions}

  Since $T$ extends the canonical Lie algebra bundle structure of $\graded(A^0)$, and
  because $\oplus_{i\ge 0}{\mathfrak a}_i$ is a grading of the canonical Lie algebra
  bundle structure of $A^0$, it is easy to see that $\tau$ is model data. Since $T$ is
  compatible with the grading of $\graded(A)$, $\tau $ is compatible with the grading
  $\oplus_i {\mathfrak a}_i$, and so compatible with the filtration of $A$. This proves
  \eqref{re1}. Clearly, $\graded_0(\tau) = T$.

  The intertwining \eqref{intertwining} means that $\tau $ is $\rho^B$-invariant,
  because $T$ is $R$-invariant. We additionally claim that
  \begin{equation}
    \rho^B_XY = [X, Y]; \qquad X \in \Gamma(B),\, Y \in \Gamma(A^0).\mathlab{cooper}
  \end{equation}
  That is, acting on $A^0$, the representation $\rho^B$ is the canonical
  representation. To prove \eqref{cooper}, it suffices to prove it in the case
  $Y \in \Gamma({\mathfrak a}_j)$, where $j\ge 0$ is arbitrary. In that case, Proposition
  \eqrefs{weyl}{pq3} guarantees that $[X, Y] \in {\mathfrak a}_j$, so that
  \begin{equation}
    \Phi([X, Y]) = [X, Y] \modulo A^{j+1}.\label{jja}
  \end{equation}
  On the other hand, we have $\Phi(Y) = Y \modulo A^{j+1}$, so that the intertwining gives
  us $\Phi(\rho^B_XY) = R_{X\modulo A^1}(Y \modulo A^{j+1})$. Since $R$
  acts on $\graded(A)_j \subset \graded(A^0)$ via the canonical action of $A/A^1$,
  we conclude
  \begin{equation}
    \label{jjb}
    \Phi(\rho^B_XY) = [X, Y] \modulo A^{j+1}.
  \end{equation}
  From equations \eqref{jja} and \eqref{jjb} and the injectivity of $\Phi $, the identity
  \eqref{cooper} is established.

  Since $A = B \oplus A^1$ (Theorem \eqrefs{weyl}{pq2}) an $A$-connection $\rho$ on $A$ is
  well-defined by
  \begin{equation*}
    \rho_{X + Y}Z = \rho^B_XZ + \tau(Y,Z);
    \qquad X \in \Gamma(B),\,Y \in \Gamma(A^1),\, Z \in \Gamma(A).
  \end{equation*}
  To prove \eqref{re2}, it suffices to show that $\rho$ is a representation, and that this
  representation extends the canonical representation of $A$ on the anchor kernel, which
  is $A^0$, by assumption. For then there exists a unique Cartan connection $\nabla $ such
  that $\bar\nabla =\rho $.

  The restriction of $\rho$ to a $A^1$-connection is a representation because $\tau $ is a
  fibre-wise Lie bracket and $A^1$ is totally intransitive. The restriction of $\rho$ to a
  $B$-connection is $\rho^B$. Therefore, to show that $\rho$ is a representation, it
  suffices to show that
  \begin{equation*}
    \rho_X \rho_Y Z - \rho_Y \rho_X Z - \rho_{[X,Y]}Z = 0,
  \end{equation*}
  whenever $X \in \Gamma(B)$ and $Y \in \Gamma(A^1)$. But in that case,
  $[X, Y] \in \Gamma(A^1)$ (because $A$ is firmly filtered) and we compute
  \begin{align*}
    \rho_X \rho_Y Z - \rho_Y \rho_X Z - \rho_{[X,Y]}Z
       &= \rho^B_X(\tau(Y, Z)) - \tau(Y, \rho^B_XZ) - \tau([X,Y], Z)\\
                     &= \rho^B_X(\tau(Y, Z)) - \tau(Y, \rho^B_XZ) - \tau(\rho^B_XY, Z),
                       \quad\text{by \eqref{cooper}}  \\
                     &= 0,
                       \quad\text{by the $\rho^B$-invariance of $\tau $.}
  \end{align*}
  The fact that $\rho$ extends the canonical representation of $A$ on $A^0$ follows from
  \eqref{cooper} and the fact that $\tau $, restricted to $A^0$, is the Lie algebroid
  bracket. This completes the proof of \eqref{re2}.

  The proof of \eqref{re3} is straightforward and left to the reader.
\end{proof}

\subsection{Reconstruction in the Lie groupoid case}
With the crucial assistance of Lemma \ref{splittings}, it is possible to globalize
the proof of \ref{reconstruction_algebroid} to obtain the following:
\begin{theorem}
  Let ${\mathcal G}$ be a transitive, firmly filtered Lie groupoid with Lie algebroid $A$,
  such that $A^0$ is the anchor kernel, and let $(R, T, E)$ be reduced geometric data for
  ${\mathcal G} $. Then there exists filtration-compatible model data $\tau $, a
  filtration compatible Cartan connection $D$, and a Weyl structure $e$, such that $\tau$,
  $D $ and $e$ are mututally compatible, the infinitesimalization $\nabla $ of $D$ is
  strongly compatible with $\tau $, $R^D=R$, $\graded_0(\tau)=T$, and $\graded_0(e) = E$.
  The Lie algebra bundle structures $\tau $ and $T$ are modeled on the same Lie
  algebra. If $T$ is regular, then so is $\tau $.
\end{theorem}
\begin{proof}
  Construct $e$, the splitting $A = \oplus_i {\mathfrak a}_i$, $\Phi$ and $\tau$ as in the
  proof of Theorem \ref{reconstruction_algebroid}. Arguments already given there establish
  that $\tau$ is filtration-compatible model data compatible with $e$. Let
  ${\mathcal H} \subset {\mathcal G}^1$ be the isotropy of $e$. Then, by Theorem
  \eqrefs{weyl2}{pr2}, the projection
  ${\mathcal G} \rightarrow {\mathcal G}/{\mathcal G}^1$ restricts to an isomophism
  ${\mathcal H} \cong {\mathcal G}/{\mathcal G}^1$. The isomorphism
  $\Phi \colon A \rightarrow \graded(A)$ interwines some representation
  $\rho^{\mathcal H} $ of ${\mathcal H} $ on $A$ with the representation $R$ of
  ${\mathcal G}/{\mathcal G}^1$ on $\graded(A)$. For some $h \in {\mathcal H} $ beginning
  at $m \in M$, we can express the intertwining as
  \begin{equation}
    \Phi(\rho^{\mathcal H}_hZ) = R_{h {\mathcal G}^1}\,\Phi(Z);
    \qquad h \in {\mathcal H},\, Z\in A|_m.\label{intertwining2}
  \end{equation}

  We claim that, restricted to $A^0$, $\rho^{\mathcal H} $ is the canonical
  representation. That is, for $h$ as above, we have
  \begin{equation}
    \label{pony4}
    \rho^{\mathcal H}_hY = \Adjoint_h Y \qquad \text{for all $Y \in A^0|_m$}.
  \end{equation}
  It suffices to prove this in the special case that $Y \in {\mathfrak a}_j$, for some
  $j \ge 0$. But then $\Adjoint_hY \in {\mathfrak a}_j$, by Theorem \eqrefs{weyl2}{pr3},
  so that
  \begin{equation}
    \label{pony9}
    \Phi(\Adjoint_hY) = \Adjoint_hY \modulo A^{j+1}.
  \end{equation}
  On the other hand, we have $\Phi(Y) = Y \modulo A^{j+1}$, so that the intertwining
  \eqref{intertwining2} and property \eqrefs{reduced_geometric}{rgd5} give us
  \begin{equation}
    \label{pony10}
    \Phi(\rho^{\mathcal H}_hY) = R_{h {\mathcal G}^1}(Y \modulo A^{j+1}) = \Adjoint_hY \modulo A^{j+1}.
  \end{equation}
  From equations \eqref{pony9} and \eqref{pony10}, the identity \eqref{pony4} is
  established.

  By definition, ${\mathcal G}^1$ (a bundle of Lie groups) has connected (source = target)
  fibres. Accordingly, Lemma \ref{splittings} shows that ${\mathcal G}^1$ must have
  {\em simply} connected fibres. In particular, the Lie algebroid representation $\mu $ of
  $A^1$ on $A$, defined by $\mu_XY = \tau(X, Y)$, integrates to some representation
  $\rho^{{\mathcal G}^1}$ of ${\mathcal G}^1$ on $A$.

  The following formula defines a candidate for a representation $\rho$ of
  ${\mathcal G} = {\mathcal G}^1 {\mathcal H} $ on $A$:
  \begin{equation*}
    \rho_{kh}Z = \rho^{{\mathcal G}^1}_k \rho^{\mathcal H}_h Z.
  \end{equation*}
  Here $h \in H$ is an arrow beginning at some $m \in M$, $k \in {\mathcal G}^1|_{hm}$ and
  $Z \in A|_m$; $\rho $ is well-defined on account of Theorem \eqrefs{weyl2}{pr2}. Since
  $\rho^{\mathcal H} $ and $\rho^{{\mathcal G}^1}$ are already representations, to show
  $\rho $ is a bone fide representation, it suffices to show that
  \begin{equation*}
    \rho^{\mathcal H}_h \rho^{{\mathcal G}^1}_k \rho^{\mathcal H}_{h^{-1}}Z = \rho^{{\mathcal G}^1}_{hkh^{-1}}Z,
  \end{equation*}
  where now $k \in {\mathcal G}^1|_m$ and $Z \in A|_{hm}$. In fact, since ${\mathcal G}^1$
  has connected fibres, it is enough to show the two Lie group homomorphisms
  ${\mathcal G}^1|_m \rightarrow \automorphism(A|_m)$ below have the same derivative:
  \begin{align*}
    k &\mapsto  \rho^{\mathcal H}_h \rho^{{\mathcal G}^1}_k \rho^{\mathcal H}_{h^{-1}}\\
    k &\mapsto  \rho^{{\mathcal G}^1}_{hkh^{-1}}.
  \end{align*}
  The readily computed derivatives are:
  \begin{align*}
    X &\mapsto \rho^{\mathcal H}_h \mu_X \rho^{\mathcal H}_{h^{-1}}(\,\cdot\,)
        = \rho^{\mathcal H}_h \tau(X, \rho^{\mathcal H}_{h^{-1}}(\,\cdot\,))\\
    \text{and}\quad X &\mapsto \mu_{\Adjoint_hX}(\,\cdot\,)=\tau(\Adjoint_hX, \,\cdot\,) =
      \tau(\rho^{\mathcal H}_h X, \,\cdot\,),\quad\text{by \eqref{pony4}}.
  \end{align*}
  Equality follows from the $\rho^{\mathcal H}$-invariance of $\tau $, which follows from
  the $R$-invariance of $T$.

  With the help of \eqref{pony4} it's not hard to see that $\rho $ extends the canonical
  representation of ${\mathcal G} $ on $A^0$. It follows that there exists a unique Cartan
  connection $D$ on $\mathcal G$ such that $\bar D = \rho $, and that $D$, $\tau $ and $e$
  have the properties stated in the theorem.

\end{proof}

\section{Parabolic geometries}
\subsection{Filtrations associated with a parabolic subalgebra}\label{canonical_filtrations}
Given any Lie algebra ${\mathfrak g} $, a subalgebra ${\mathfrak g}^0$, and an ideal
${\mathfrak g} ^1$ of ${\mathfrak g}^0$, we obtain a Lie algebra filtration of
$\mathfrak g$ by inductively defining
${\mathfrak g}^j = [{\mathfrak g}^1, {\mathfrak g}^{j-1}]$, for $j >1$, and
${\mathfrak g}^k = \{X \in {\mathfrak g} \suchthat [X, {\mathfrak g}^1] \subset {\mathfrak
  g}^{k+1}\}$, for $k<0$. See, e.g., \cite{Calderbank_Noppakaew_16}. If ${\mathfrak g} $
is a semisimple Lie algebra over ${\mathbb C} $ or ${\mathbb R} $, then
${\mathfrak p} \subset {\mathfrak g} $ is {\df parabolic} if it contains a Borel subagebra
of ${\mathfrak g} $ or, equivalently, if ${\mathfrak p}^\perp$ is the nilpotent radical of
${\mathfrak p}$ (see op.~cited). Here $\perp$ denotes orthogonalization with respect to
the Killing form. In this case, taking ${\mathfrak g}^0={\mathfrak p} $ and
${\mathfrak g}^1 = {\mathfrak p}^\perp$, the inductively defined filtration just
desrcribed is called the {\em canonical filtration} defined by ${\mathfrak p}$. The
following fact is known:

\begin{lemma}
  If ${\mathfrak g} $ is a semisimple Lie algebra and
  ${\mathfrak p} \subset {\mathfrak g} $ a parabolic subalgebra, then the graded Lie
  algebra $\graded({\mathfrak g})$ associated with the canonical filtration is
  (non-canonically) isomorphic to ${\mathfrak g} $.
\end{lemma}

\begin{proof}[Sketch of proof]
  From the definition of the canonical filtration, it is not hard to show that
  $({\mathfrak g}^j)^\perp = {\mathfrak g}^{1-j}$, $j \in {\mathbb Z} $. This condition
  allows one to drop the Killing form to a non-degenerate bilinear form on
  $\graded({\mathfrak g})$ in which
  $\graded({\mathfrak g})_j \perp \graded({\mathfrak g})_k$, unless $j + k =0$. One shows
  that this must coincide with the Killing form on $\graded({\mathfrak g})$, which means
  $\graded({\mathfrak g})$ is semisimple, by Cartan's criterion. It follows that
  $\graded({\mathfrak g})$ has a (unique) grading element, and consequently that
  ${\mathfrak g} $ and $\graded({\mathfrak g})$ are isomorhic, by the argument given in
  \ref{splittings}. For details, see the more general argument at \cite[Propositions 2.14
  \& 3.9]{Calderbank_Noppakaew_16}.
\end{proof}

\subsection{Parabolic geometries defined}\label{parabolic}
Suppose that $G$ is a semisimple Lie group and $P \subset G$ a parabolic subgroup. By
definition, this means $P$ has a parabolic Lie algebra
${\mathfrak p} \subset {\mathfrak g} $ and that the associated canonical filtration of
${\mathfrak g} $ is $P$-invariant under $\Adjoint$ (it is automatically
${\mathfrak p} $-invariant under $\adjoint$). The quotient $M=G/P$ is known as a {\df
  generalized flag manifold}. The $P$-invariance of the canonical filtration leads to a
firm filtering of ${\mathcal G} = G \times M$, as described in \ref{homogeneous}, where we
noted the existence of a natural Cartan connection $D$ and model data $\tau$. We now turn
attention to curved generalizations of these spaces.

A {\df parabolic geometry} consists of a transitive, firmly
filtered Lie groupoid ${\mathcal G} $ over some smooth manifold $M$, with associated Lie
algebroid $A$, and filtration-compatible model data $\tau $, such that:
\begin{conditions}
    \item $A^0$ is the kernel of the anchor of $A$.\label{parabolic1}
    \item The Lie algebra model for $\tau $ is semisimple.\label{parabolic2}
    \item $A^0$ is a fibre-wise parabolic subalgebra of $(A, \tau)$.\label{parabolic3}
    \item There exists a Cartan connection $D$ on ${\mathcal G} $
      compatible with $\tau $ and the filtration.\label{parabolic4}
\end{conditions}

An {\df infinitesimal parabolic geometry} consists of a transitive, firmly filtered Lie
algebroid $A$, and filtration-compatible model data $\tau $, satisfying the above
conditions, with the Cartan connection $D$ in \eqref{parabolic4} replaced with an
infinitesimal Cartan connection $\bar\nabla$. Evidently, every parabolic geometry has a
corresponding infinitesimalization.

A parabolic geometrty is {\df regular} if the associated model data $\tau $ is regular
(see \ref{regularity}).

\begin{proposition}
  If $M$ is connected then, in the global groupoid case, and in the case of an integrable
  Lie algebroid, it suffices to check that \eqref{parabolic1} and \eqref{parabolic3} hold
  over a single point $m \in M$.
\end{proposition}

\begin{theorem}
  In a global or infinitesimal parabolic geometry, the compatible Cartan connection is
  uniquely determined by the model data $\tau$, with which it is strongly compatible.
\end{theorem}

\begin{proof}
  In the infinitesimal case, suppose that $\nabla^1$ and $\nabla^2$ are Cartan connections
  compatible with model data $\tau$ on $A$, and write
  $\delta_XY = \bar\nabla^1_XY - \bar\nabla^2_XY$. For each $X \in A|_m$, $m \in M$, we
  may regard $\delta_X$ as a linear map $A|_m \rightarrow A|_m$ and, because $\tau $ is
  both $\bar\nabla^1$ and $\bar\nabla^2$-invariant (by compatability) it follows that
  $\delta_X$ is even an Lie algebra endomorphism, with respect to the bracket $\tau
  $. Moreover, as $\bar\nabla^j$ extends the canonical representation of $A$ on its anchor
  kernel $A^0$, it follows that $\delta_XY = 0$ for all $Y$ in the parabolic subalgebra
  $A^0|_m \subset A|_m$. By semisimplicity, there exists $Z \in A|_m$ such that
  $\delta_XY = [Z, Y]$, and in particular such that $[Z, Y]=0$ for all $Z \in A^0|_m$. So
  $Z$ is in the centralizer of a parabolic subalgebra of a semisimple Lie algebra, and so
  is in the centralizer of some Borel subalgebra, and consequently zero. This shows
  $\delta_X=0$ and we conclude $\bar\nabla^1 = \bar\nabla^2$. Since $A$ is transtive,
  $\nabla^1=\nabla^2$.

  To prove strong compatability, suppose $\bar\nabla $ is a $\tau $-compatible Cartan
  connection (the infinitesmialization of some Cartan connection $D$ in the global
  case). Fix some $X \in A^0|_m$, $m \in M$.  Then as $A^0$ lies in the anchor kernel, we
  may view $\bar\nabla_X$ as a linear map $A|_m \rightarrow A|_m$. In fact, with respect
  to the bracket $\tau $, this is a Lie algebra endomorphism, because $\tau $ is
  $\bar\nabla $-invariant ($\nabla $ and $\tau $ are compatible). By the semisimplicity,
  there is $X' \in A|_m$ such that $\bar\nabla_XY=\tau(X',Y)$ for all $Y \in A|_m$.
  % Moreover, since $A^0|_m \subset A|_m$ is a {\em parabolic} subalgebra, and
  % $\bar\nabla $ is filtration-preserving (Proposition \eqrefs{model_data_gradings}{u7}) we
  % must have $X' \in A^0|_m$.
  In the special case $Y \in A^0|_m$, we have $\bar\nabla_XY = [X, Y] = \tau(X, Y)$, the
  first equality following from \eqrefs{hsd}{zero}, the second because $\tau $ extends the
  canonical Lie algebra structure of $A^0$. Therefore, $\tau(X' - X, Y) = 0$ for all such
  $Y$. In other words, $X' - X$ must lie in the centralizer of a parabolic subalgebra,
  which is trivial, as noted above. We conclude $X'=X$ and whence
  $\bar\nabla_XY = \tau(X, Y)$ for all $Y \in A|_m$.

  In the global case, which is proven similarly, suppose $ D^1$ and $D^2$ are Cartan
  connections compatible with model data $\tau $ and write
  $\Delta_gY = \bar D^1_{g^{-1}} \bar D^2_g Y$. For each arrow $g \in {\mathcal G} $
  beginning at some $m \in M$, we may regard $\Delta_g$ as a vector space automorphism
  $A|_m \rightarrow A|_m$ and, because $\tau $ is both $\bar D^1$ and
  $\bar D^2$-invariant, it follows that $\Delta_g$ is even a Lie algebra automorphism,
  with respect to the bracket $\tau$. Since $\bar D^j$ extends the canonical representation of
  $\mathcal G$ on $A^0$, $\Delta_gY = Y$ for all $Y \in A^0|_m$. So $\Delta_g$ is an
  automorphism of a semisimple Lie algebra fixing points in a parabolic subalgebra, and is
  consequently the identity. (In the complex Lie algebra case, this already follows from
  the fact that all elements of some Cartan subalgebra are being fixed, and in the real
  case, because the same holds for the complexification of $\Delta_g$.) We conclude
  $\bar D^1= \bar D^2$ and, by transitivity, $D^1 = D^2$.
\end{proof}

\subsection{All parabolic geometries admit Weyl structures}\label{all}
Suppose $\tau $ is the model data associated with a parabolic geometry (global or
infinitesimal). Then, according to Lemma \ref{canonical_filtrations},
$T = \graded_0(\tau)$ is a semisimple Lie algebra bundle structure on $\graded(A)$, where
$A$ is the underlying filtered Lie algebroid. We accordingly have a {\em unique} section
$E$ of $\graded(A)_0$ that is a fibre-wise grading section of $(\graded(A), T)$, what we
shall call the {\df canonical grading section}.
\begin{proposition}
  For any global or infinitesimal parabolic geometry, any lift $e$ of the canonical
  grading section $E$ is a Weyl structure compatible with the model data.
\end{proposition}

\begin{proof}
  This time we begin with the global case which is simpler.  Let $({\mathcal G}, \tau)$ be
  a parabolic geometry, $M$ the base manifold, and $D$ a Cartan connection compatible with
  $\tau $ and the filtration of the Lie algebroid $A$ of $\mathcal G$.  Since $\tau $ is
  $\bar D$-invariant (by compatiblity) $T$ is invariant with respect to the representation
  $R^D$ of ${\mathcal G} $ defined in \ref{reduced_geometric}. In particular, for any
  arrow $g \in {\mathcal G} $ beginning at some $m \in M$, $R^D_g$ must map grading
  elements of $\graded(A)|_m$ to grading elements of $\graded(A)|_{gm}$, which, by the
  forementioned uniqeness, means $R^D_gE(m) = E(gm)$, i.e., $E$ is $R^D$-invariant. So
  $(R^D, T, E)$ is reduced geometric data and any lift $e$ of $E$ through
  $A \rightarrow A/A^1$ will be a Weyl structure for ${\mathcal G} $, by Proposition
  \ref{reduced_geometric}. By construction, $e$ will be compatible with $\tau $.

  Turning to the infinitesimal case, let $(A, \tau)$ be an infinitesimal parabolic
  geometry and $\nabla$ a Cartan connection compatible with $\tau $ and the filtration of
  $A$. Since $\tau $ is $\bar\nabla $-invariant (by compatibility) $T$ is invariant with
  respect to the representation $R^\nabla $ of $A$ on $\graded(A)$ defined in
  \ref{reduced_geometric}. Let $X \in \Gamma(A)$, $j \in {\mathbb Z} $ and
  $Y \in \Gamma(\graded(A)_j)$ be arbitrary. Then $T(E, Y) = jY$ and, applying
  $R_X^\nabla $ to both sides, we obtain, courtesy of the $R^\nabla$-invariance of $T$,
  \begin{equation}
    \label{won}
    T(R_X^\nabla E, Y) + T(E, R_X^\nabla Y) = j R_X^\nabla Y.
  \end{equation}
  On the other hand, the grading of $\graded(A)$ is $R^\nabla $-invariant, so that
  $R_X^\nabla Y \in \Gamma(\graded(A)_j)$, implying $T(E, R_X^\nabla Y) = j R_X^\nabla
  Y$. From \eqref{won} we now conclude that $T( R_X^\nabla E, Y) =0$. From the
  arbitrariness of our choice of $j$ and $Y$ above, we conclude that $ R_X^\nabla E$ is
  fibre-wise in the center of a semisimple Lie algebra, and consequently vanishes. This
  shows $(R^\nabla, T, E)$ is reduced geometric data and the result follows Proposition
  \ref{reduced_geometric}.
\end{proof}

\subsection{Model curvature}\label{model_curvature}
Given a global or infinitesimal parabolic geometry with underlying Lie algebroid
$\# \colon A \rightarrow TM$, let $\tau $ denote the model data and $\nabla $ the
compatible infinitesimal Cartan connection. With Theorem \ref{curvature} in mind, the
difference between between $\torsion \bar\nabla$ (a ``deformed bracket'') and $\tau $ (the
model bracket) will be the {\df model curvature} of the parabolic geometry. Specifically,
this is the section $\kappa$ of $\wedge^2 T^*\!M \otimes A$ defined by
\begin{equation*}
  \kappa(\#X, \#Y) = \torsion \bar\nabla (X, Y) - \tau(X, Y); \qquad X, Y \in \Gamma(A),
\end{equation*}
which makes sense, on account of the strong compatability property \eqref{strong} and
equation \eqref{zero} defining $\bar\nabla $.

The model curvature is the Lie algebroid analogue of ``curvature'', as that term is
generally understood in the Cartan geometry literature.

\begin{proposition}
  If $\kappa =0$, then $\curvature \nabla =0$. If the geometry is regular, then the model
  curvature is homogeneous of degree one, i.e., $\kappa(U, V) \in \Gamma(A^{j+k+1})$ for
  $U \in \Gamma(T^jM)$ and $V \in \Gamma(T^kM)$; it is otherwise of degree zero.
\end{proposition}
\noindent The consequences of $\curvature \nabla =0$ have already been discussed in
\ref{curvature}. At the very least, according to Theorem \ref{curvature}, $A$ will locally
be an action algebroid, the operative Lie algebra being the semisimple Lie algebra model
for $\tau $. In any case, we are now justified in describing parabolic geometries as
curved generalizations of generalized flag manifolds.

\begin{proof}
  For any Cartan connection $\nabla $ on a transitive Lie algebroid, one has the easily
  derived identity,
  \begin{equation*}
    \curvature\nabla(\#X, \#Y)Z =  \bar\nabla_Z\torsion \bar\nabla (X, Y).
  \end{equation*}
  See, for example, \cite[\S5.4]{Blaom_06}. If $\kappa =0$, i.e.,
  $\torsion \bar\nabla = \tau $. Then the $\bar\nabla $-invariance of $\tau $
  (compatability of $\tau$ and $\nabla$) means $\bar\nabla\torsion \bar\nabla =0$, and
  $\curvature \nabla =0$ follows from the identity.

  By definition,
  \begin{equation}
    \label{pot}
    \kappa(\#X, \#Y) = \bar\nabla_XY - \bar\nabla_YX - ([X, Y] + \tau(X, Y))
  \end{equation}
  Suppose $X \in \Gamma(A^j)$ and $Y \in \Gamma(A^k)$ for some $j, k < 0$. As we argued in
  the proof of Proposition \ref{homogeneous}, the non-positivity of $j$ and $k$ means that
  $A^j \subset A^{j+k+1}$ and $A^k \subset A^{j+k+1}$. By $\bar\nabla $-invariance of the
  filtration of $A$, the first two terms on the right-hand side of \eqref{pot} lie in
  $\Gamma(A^{j+k+1})$. The last bracketted term lies in $\Gamma(A^{j+k})$, by axiom F1 for
  Lie algebroid filtrations, and compatability of $\tau $ with the filtration of $A$. In
  the regular case, this bracketted term lies in $A^{j+k+1}$, by the definition of
  regularity.

\end{proof}

\subsection{Semidirect product structure}\label{semidirect2}
Let $({\mathcal G} , \tau)$ be a parabolic geometry over $M$, $A$ the underlying filtered
Lie algebroid, and $D$ the compatible Cartan connection.  Transitivity means we can
identify $\oplus_{i<0}\graded(A)$ with $\graded(TM)$, so that the grading-preserving
representation $R^D $ of ${\mathcal G}/{\mathcal G}^1$ on $\graded(A)$, which also
preserves $\graded_0(\tau)$, delivers, in particular, a representation of the same on
$\graded(TM)$.

\begin{proposition}
  In the scenario above, and in the analogous infinitesimal case, identify
  $\graded(A^1)=\oplus_{i>0}\graded(A)_i$ with $\graded(TM)^*\cong\graded(T^*\!M)$ using
  the Killing form defined by $T=\graded_0(\tau)$. Then, under this identification, the
  canonical representation of ${\mathcal G}/{\mathcal G}^1$ (respectively, $A/A^1$) on
  $\graded(A^1)$ is the dual of the representation of the representation $R^D$ on
  $\graded(TM)$ (respectively $R^\nabla$ on $\graded(TM))$ noted above.
\end{proposition}

\begin{proof}
  This follows easily from the $R^D$-invariance (respectively, $R^\nabla$-invariance) of
  $T$, and hence the associated Killing form, and the fact that $R^D$ (respectively,
  $R^\nabla$) coincides with the canonical representation on $\graded(A^1)$.
\end{proof}
\noindent Applying Theorem \ref{weyl}:

\begin{corollary}
  The filtered Lie algebroid $A$ in a parabolic geometry $(A, \tau)$ is isomorphic to the
  semidirect product $A/A^1 \oplus_{R^\nabla}\graded(T^*\!M)$, where $\nabla $ is the
  compatible Cartan connection.
\end{corollary}

An analagous conclusion may be formulated in the global case.

\subsection{Levi geometries and their prolongation}
Let $(A, \tau)$ be an infinitesimal parabolic geometry. Recall from \ref{reduction} that
$B = A/A^1$ is filtered with $B^j=0$ for $j \ge 1$. In the present case, $B$ is transitive
and $B^0=\graded(A)_0$ is the anchor kernel of $B$. The fibre-wise Levi decomposition of
the fibre-wise parabolic subalgebra $\graded(A^0) \subset \graded(A)$ is
$\graded(A^0)=\graded(A)_0 \oplus\graded(A^1)=B^0 \oplus \graded(A^1)$, where $B^0$ is the
fibre-wise reductive part. In other words, the Lie algebra bundle $B^0$ is modelled on a
Lie algebra constituting the reductive component in the Levi decomposition of a parabolic
subalgebra, which is usually called the {\df Levi subalgebra}. Motivated by our
observations above, we introduce the following terminology:

\begin{definition}
  An {\df infinitesimal Levi geometry} consists of a transitive filtered Lie algebroid $B$
  over $M$, with $B^j=0$ for $j\ge 1$, together with a representation $R$ of $B$ on
  $\graded(TM)\cong \oplus_{i<0}\graded(B)_i $, and an $R$-invariant semisimple Lie
  algebra bundle structure $T$ on
  \begin{equation}
    \graded(TM) \oplus B^0 \oplus \graded(T^*\!M)\label{fug}
  \end{equation}
  compatible with the obvious grading and duality (the canonical pairing of $\graded(TM)$
  and $\graded(T^*\!M)$ should coincide with that defined by the Killing form).

  A {\df Levi geometry} consists of a filtered Lie groupoid ${\mathcal H} $ over $M$,
  whose infinitesimalization $B$ is filtered as just described, a representation $R$ of
  ${\mathcal H} $ on $\graded(TM)$, together with an $R$-invariant, Lie algebra bundle
  structure $T$ on \eqref{fug} compatible with the grading and duality.

  An (infinitesimal)
  Levi geometry is {\df regular} if the canonical Lie algebra bundle structure of
  $\graded(TM)$ is the negative of the restiction of $T$.
\end{definition}

\begin{example}[Conformal structures]
  Define $B \subset J^1 (TM) $ and $T$ as in \ref{conformal} and filter $B$ according to
  $B^j = B$ for $j<0$, $B^0=\co(TM)$, and $B^j=0$ otherwise. Let $B$ act on
  $\graded(TM)=TM$ by restricted adjoint action. Then this makes $B$ into a Levi geometry.

  If we let ${\mathcal H} \subset J^1 (M \times M)$ denote the Lie groupoid of one-jets of
  local diffeomorphisms preserving the conformal structure to first order, and let
  ${\mathcal H} $ act on $\graded(TM)=TM$ by restricted adjoint action, then
  ${\mathcal H} $ becomes a global Levi geometry.
\end{example}

Unfortunately for us, the analogue of a {\df global} Levi geometry appears in the
{\v{C}}ap-Slov{\'a}k-Schichl theory under the name ``infinitesimal flag structure''
\cite{Cap_Slovak_09}.  In our context this is, moreover, a label more apt for the
parabolic specialisations of the action algebroids described in \ref{eg2}, a further
argument for not adopting the terminology here.

According to our observations at the beginning of \ref{semidirect2}, we can associate with
each parabolic geometry $({\mathcal G}, \tau)$, a Levi geometry
$(({\mathcal G}/{\mathcal G}^1), R^D, \graded_0(\tau))$, which we call its {\df Levi
  reduction}, and analogous observations hold in the infinitesimal case. The anchor kernel
$B^0$ in a Levi geometry is always a reductive Lie algebra bundle. If $\tau $ is regular,
then so is the Levi reduction.

For a converse construction, consider an infinitesimal Levi geometry $(B, R, T)$. The
bracket $T$ makes $\graded(T^*\!M)$ into a Lie algebra bundle, on which $B$ acts via $R$,
and we call the semidirect product A := $B \oplus_R \graded(T^*\!M)$ the Levi geometry's
{\df algebraic prolongation}. It is a transitive Lie algebroid, and filtered as we
described at the end of \ref{reduction}, a filtration that is obviously firm. Just as
obvious is that $B \cong A/A^1$,
\begin{equation*}
  \graded(A) \cong \graded(TM) \oplus B^0 \oplus \graded(T^*\!M),
\end{equation*}
and that $(T, R, E)$ is reduced geometric data for $A$, if we take $E$ to be the unique
grading section of $\graded(A)$, viewed as a Lie algebra bundle using the semisimple
bracket $T$. Applying Theorem \ref{reconstruction_algebroid}, we can equip $A$ with model
data $\tau $, with the same Lie algebra model as $T$, compatible with the filtration, and
strongly compatible with some Cartan connection $\nabla $, such that $R^\nabla=R$ and
$\graded_0(\tau)=T$.  That is, $(A, \tau)$ is an infinitesimal parabolic geometry whose
Levi reduction is $(B, R, T)$. Corollary \ref{semidirect2} ensures that the algebraic
prolongation $A$ is, up to isomorphism, the {\em unique} Lie algebroid supporting the
structure of an infinitesimal parabolic structure with this property. Theorem
\ref{reconstruction_algebroid} additionally ensures that if the Levi geometry is regular,
then the reconstructed parabolic geometry will be regular also. We have now proven the
infinitesimal part of the following result:

\begin{theorem}  Every  (infinitesimal)  Levi  geometry  is the  Levi  reduction  of  some
  (infinitesimal)  parabolic geometry.  Up  to isomorphism,  the  underlying Lie  groupoid
  (algebroid) of  the parabolic geometry is  determined uniquely by the  Levi geometry. If
  the Levi reduction is regular, then we  may take the corresponding parabolic geometry to
  be regular also.
\end{theorem}

\begin{proof}[Proof (global case)]

\end{proof}
\section{Deforming Cartan connections}

Associated with any Lie algebroid $A$ over $M$ is a bundle of Lie groups we shall denote by
$\automorphism (TM; A)$ that (in the transitive ?) is a commutative extension of
$\automorphism(TM) = \automorphism(TM; TM)$. Our interest in this bundle lies in the fact
that, for any Lie groupoid ${\mathcal G} $ integrating $A$ it is a model for the kernel of
the projection $J^1 {\mathcal G} \rightarrow {\mathcal G} $ \cite{}. As we explain in
Appendix \ref{}, sections of $\automorphism(TM; A)$ lead, via the adjoint action of
$J^1 {\mathcal G} $ on ${\mathcal G} $, to deformations of Cartan connections on
${\mathcal G} $ (left inverses for $J^1 {\mathcal G} \rightarrow {\mathcal G} $). Here we
shall describe the corresponding deformations of {\em infinitesimal} Cartan connections
$\nabla $ on $A$, which we can do without reference to an integrating Lie groupoid
${\mathcal G} $.

For each $\phi \in T^*\!M \otimes {\mathfrak g} $ with base-point $m
\in M$, define $\phi_{TM}\in T^*\!M \otimes {TM}$ by
\begin{equation*}
  \phi_{TM} v = v + \#\,\phi\,v.
\end{equation*}
Then, $\model$ is defined by
\begin{equation*}
  \model := \{\phi \in T^*\!M \otimes
  {A} \suchthat \phi_{TM} \in \automorphism({TM})\}.
\end{equation*}
This is a totally intransitive Lie groupoid (bundle of Lie groups) with the zero-section
as set of identity elements. Multiplication is given by
\begin{equation}
  \psi \phi = \psi + \phi + \psi \circ \# \circ \phi = \phi +
  \psi \circ \phi_{TM},\mathlab{clost}.
\end{equation}
Inversion is given by
\begin{equation*}
  \phi \mapsto  -\phi \circ (\phi_{TM})^{-1}.
\end{equation*}
With respect to this extra structure, the map
$$\phi \mapsto \phi_{TM} \colon \model \rightarrow
\automorphism({TM})$$ is a morphism (of totally intransitive
groupoids) and the commutative kernel consists of those
$\phi \colon TM \rightarrow A$ with image lying in the kernel of the
anchor.

In addition to the natural representation of
$\model $ on $TM$ given by $\phi \cdot v = \phi_{TM} v$, we have a
representation of $\model$ on ${A} $ given by
\begin{equation*}
  \Adjoint_\phi X := X + \phi\,\#\,X; \qquad X  \in {A}.
\end{equation*}
With respect to these actions of $\model$ on ${A} $ and $TM$, the
anchor $\#$ is equivariant. The restricted action of $\model$ on the
kernel of $\#$ is trivial.

Given an infinitesimal Cartan connection $\nabla$ on $A$, its {\df
$\phi$-deformation} is defined by
\begin{equation*}
  \nabla'_VX = \nabla_VX - (\bar\nabla\phi)(\phi_{TM}^{-1}V).
\end{equation*}
This formula becomes a little less mysterious if we consider the
corresponding deformation to the representation $\bar\nabla$ defined
by $\nabla $. It is given by
\begin{align}
  \bar\nabla'_XY &= \Adjoint_\phi (\bar\nabla_X(\Adjoint_\phi^{-1}Y)),\notag\\
  &=\bar\nabla_XY - (\bar\nabla_{?X}\phi)(\phi_{TM}^{-1}(\#Y)),
\end{align}
as the reader is invited to check herself.



\appendix
\section{Details}
\subsection{Proof of Lemma \ref{splittings}} \label{appendix_splittings} We note that
${\mathfrak g}^1$ is nilpotent, from the definition of a grading. Moreover, the grading
property ensures that the action of ${\mathfrak g}^1$ on ${\mathfrak g}^0 $ preserves the
grading, and the corresponding action on $\graded({\mathfrak g}^0)$ is trivial. The same
must be true for adjoint action of $G^1$, as it is connected. In particular
$e + {\mathfrak g}^1$ must be $G^1$-invariant, for any lift $e \in {\mathfrak g}^0$ of
$E$. Suppose that the action of $G^1$ on $e + {\mathfrak g}^1$ is infinitesimally
free. Then, by a dimension count, this action is infinitesimally transitive, and so
transitive, since $G^1$ is connected. Infinitesimal freeness means the stabilizer subgroup
$H$ for the action is discrete, so that $G^1 \rightarrow G^1/H \cong e + {\mathfrak g}^1$
is a covering space. Since $e + {\mathfrak g}^1$ is simply connected, $H$ is trivial and
the remaining claims of part \eqref{splittings1} of Lemma \ref{splittings} follow.

To show that the action of $G^1$ on $e + {\mathfrak g}^1$ is infinitesimally free, for
each $X \in {\mathfrak g}^1$, let $X^\dagger$ denote the infinitesimal generator of $X$, a
vector field on $e + {\mathfrak g}^1$. We have $X^\dagger(e') = [X, e']$, for
$e' \in e + {\mathfrak g}^1$. Supposing $X^\dagger(e') = 0$, write $X=\sum_{j>0} X_j$ for
some $X_j \in {\mathfrak g}_j$, where ${\mathfrak g}=\oplus_j {\mathfrak g}_j$ is the
$E$-splitting corresponding to $e'$. Then we conclude that $\sum_{j>0}jX_j=0$, from which
it follows that $X=0$.

To prove part \eqref{splittings1} of Lemma \ref{splittings}, observe that $\Phi_e$ is
characterized among all ismorphisms $\graded({\mathfrak g}) \rightarrow {\mathfrak g}$ by
the requirement
\begin{equation*}
  [\Phi_e(X), e] = j \Phi_e(X)\quad
  \text{for all $X \in \graded_j({\mathfrak g}),\, j \in {\mathbb Z}$}.
\end{equation*}
Taking $e = e_1$ and applying $\Adjoint_g$ to both sides, we obtain
$[\Adjoint_g \phi_{e_1}X, g\cdot e_1] = jg\cdot e_1$, i.e.,
\begin{equation*}
  [(\Adjoint_g \circ \phi_{e_1})X, e_2] = je_2,
\end{equation*}
if $e_2 = g\cdot e_1$. So, by the forementioned characaterization, we conclude
$\Adjoint_g\circ \phi_{e_1} = \Phi_{e_2}$.

\bibliography{dynamics2}

\end{document}

%%% Local Variables:
%%% mode: latex
%%% TeX-master: t
%%% End:
